% Generated by the blueprint tool (blueprint-pdf/1) from a blueprint snapshot; edit the
% Markdown sources, not this file.  Compile with `latexmk -pdf` or by running
% `pdflatex` twice (the second run settles the contents and the page references).
% Images, if any, are read from assets/ next to this file.
\documentclass[11pt]{article}
\usepackage[T1]{fontenc}
\usepackage[utf8]{inputenc}
\usepackage{lmodern}
\usepackage{textcomp}
\usepackage[a4paper,margin=2.5cm,headheight=14pt]{geometry}
\usepackage{amsmath,amssymb,amsthm}
\usepackage{amscd}  % KaTeX's CD diagrams; part of amsmath, so always there
\usepackage{mathtools}
\usepackage{bm}
\usepackage{graphicx}
\usepackage{xcolor}
\usepackage{array,booktabs,tabularx}
\usepackage{fancyvrb}
\usepackage[normalem]{ulem}
\usepackage{microtype}
% Optional, and only with their fonts: a small TeX installation has the
% packages without the fonts, which fails at the first formula.
\IfFileExists{mathrsfs.sty}{\IfFileExists{rsfs10.tfm}{\usepackage{mathrsfs}}{}}{}
\IfFileExists{stmaryrd.sty}{\IfFileExists{stmary10.tfm}{\usepackage{stmaryrd}}{}}{}
\IfFileExists{cancel.sty}{\usepackage{cancel}}{}
\usepackage[hyphens]{url}
\usepackage[pdfusetitle,bookmarksnumbered=false]{hyperref}
\usepackage{bookmark}
\definecolor{bplink}{HTML}{1F4E9C}
\definecolor{bpmuted}{HTML}{5F676E}
\hypersetup{colorlinks=true,linkcolor=bplink,urlcolor=bplink,citecolor=bplink}
\setlength{\parindent}{0pt}
\setlength{\parskip}{0.55em plus 0.2em minus 0.1em}
% Lists as tight as the page's.
\makeatletter
\def\@listi{\leftmargin\leftmargini\parsep 0pt\topsep 0.3em plus 0.1em\itemsep 0.1em plus 0.1em}
\let\@listI\@listi
\def\@listii{\leftmargin\leftmarginii\labelwidth\leftmarginii\advance\labelwidth-\labelsep\parsep 0pt\topsep 0.1em\itemsep 0.1em}
% An ordered list that does not start at 1.
\newcommand{\bpstartat}[1]{\setcounter{\@enumctr}{\numexpr#1-1\relax}}
\makeatother
\setcounter{tocdepth}{3}
\emergencystretch=3em
\allowdisplaybreaks

% A numbered statement, set as a paper sets one: "Theorem 2.1 (Title)."
% #1 is the lead ("Theorem~2.1"), #2 what follows it in roman.
\newenvironment{bpstatement}[2]{\par\addvspace{0.6\baselineskip}\noindent\textbf{#1}#2\textbf{.}\enspace\ignorespaces}
  {\par\addvspace{0.4\baselineskip}}
\newenvironment{bpstatementit}[2]{\begin{bpstatement}{#1}{#2}\itshape\ignorespaces}{\end{bpstatement}}
% A step: a relation with prose of its own, read under the object it starts from.
\newenvironment{bpstep}[1]{\par\addvspace{0.3\baselineskip}\begin{list}{}{\leftmargin=1.6em\rightmargin=0pt\topsep=0pt}\item\relax
  {\small\color{bpmuted}#1}\par\nopagebreak}{\end{list}\par}
% Notes about where an object is written out.
\newenvironment{bpnote}{\par\small\color{bpmuted}\itshape}{\par}
\newcommand{\bpomitted}[1]{\par{\small\color{bpmuted}\itshape #1\par}}
\newcommand{\bpbodyheading}[1]{\par\addvspace{0.4\baselineskip}\noindent\textbf{#1}\par\nopagebreak}
\newcommand{\bprule}{\par\medskip\noindent\hfil\rule{0.3\linewidth}{0.4pt}\hfil\par\medskip}
\newcommand{\bpboundary}[1]{\par\nopagebreak{\small\color{bpmuted}#1\par}\nopagebreak}
% A reference to something this export leaves out, and one that resolves nowhere.
\newcommand{\bpout}[1]{#1\textsuperscript{\hyperlink{bp-outside}{\ensuremath{\ast}}}}
\newcommand{\bpbroken}[1]{\textcolor{red!70!black}{[#1?]}}
\newcommand{\bpslug}[1]{\texttt{#1}}
% A bibliography entry: label, anchor, entry; the label hangs in the margin.
\newcommand{\bpbibitem}[3]{\par\setbox0\hbox{[#1]\enspace}\ifdim\wd0<5em\setbox0\hbox to 5em{[#1]\hfil}\fi
  \hangindent=\wd0 \hangafter=1 \noindent\box0\phantomsection\label{#2}#3\par\smallskip}
% A character the export knows no way to set.
\newcommand{\bpglyph}[1]{\fbox{\scriptsize #1}}
% Images, at their own size unless they are wider than the line.
\newsavebox{\bpimagebox}
\newcommand{\bpimage}[1]{\sbox{\bpimagebox}{\includegraphics{#1}}%
  \ifdim\wd\bpimagebox>\linewidth\resizebox{\linewidth}{!}{\usebox{\bpimagebox}}\else\usebox{\bpimagebox}\fi}
\newcommand{\bpmissingimage}[1]{\fbox{\parbox{0.8\linewidth}{\centering\small\color{bpmuted}Image not included: #1}}}
\DefineVerbatimEnvironment{bpverbatim}{Verbatim}{fontsize=\small,xleftmargin=1em}
\renewcommand{\qedsymbol}{\ensuremath{\blacksquare}}
\hypersetup{pdftitle={Weil Conjectures — 6 Constructibility of f!},pdfsubject={Chapter 6 Constructibility of f ! ; all levels},pdfcreator={blueprint}}
% KaTeX commands the website knows
\providecommand{\R}{\mathbb{R}}\providecommand{\N}{\mathbb{N}}\providecommand{\Z}{\mathbb{Z}}
\providecommand{\Q}{\mathbb{Q}}\providecommand{\C}{\mathbb{C}}
\providecommand{\Reals}{\mathbb{R}}\providecommand{\reals}{\mathbb{R}}\providecommand{\Complex}{\mathbb{C}}
\providecommand{\cnums}{\mathbb{C}}\providecommand{\natnums}{\mathbb{N}}\providecommand{\Bbb}[1]{\mathbb{#1}}
\providecommand{\bold}[1]{\mathbf{#1}}\providecommand{\frak}[1]{\mathfrak{#1}}\providecommand{\mathbbm}[1]{\mathbb{#1}}
\providecommand{\lang}{\langle}\providecommand{\rang}{\rangle}\providecommand{\bra}[1]{\langle #1\rvert}
\providecommand{\ket}[1]{\lvert #1\rangle}\providecommand{\braket}[1]{\langle #1\rangle}
\providecommand{\set}[1]{\{#1\}}\providecommand{\Set}[1]{\left\{#1\right\}}
\providecommand{\infin}{\infty}\providecommand{\isin}{\in}\providecommand{\plusmn}{\pm}\providecommand{\sub}{\subset}
\providecommand{\sube}{\subseteq}\providecommand{\supe}{\supseteq}\providecommand{\exist}{\exists}
\providecommand{\bull}{\bullet}\providecommand{\larr}{\leftarrow}\providecommand{\rarr}{\rightarrow}
\providecommand{\lrarr}{\leftrightarrow}\providecommand{\Larr}{\Leftarrow}\providecommand{\Rarr}{\Rightarrow}
\providecommand{\Lrarr}{\Leftrightarrow}\providecommand{\uarr}{\uparrow}\providecommand{\darr}{\downarrow}
\providecommand{\harr}{\leftrightarrow}\providecommand{\lArr}{\Leftarrow}\providecommand{\rArr}{\Rightarrow}
\providecommand{\hArr}{\Leftrightarrow}\providecommand{\alef}{\aleph}\providecommand{\alefsym}{\aleph}
\providecommand{\real}{\Re}\providecommand{\image}{\Im}\providecommand{\weierp}{\wp}
\providecommand{\Alpha}{\mathrm{A}}\providecommand{\Beta}{\mathrm{B}}\providecommand{\Epsilon}{\mathrm{E}}
\providecommand{\Zeta}{\mathrm{Z}}\providecommand{\Eta}{\mathrm{H}}\providecommand{\Iota}{\mathrm{I}}
\providecommand{\Kappa}{\mathrm{K}}\providecommand{\Mu}{\mathrm{M}}\providecommand{\Nu}{\mathrm{N}}
\providecommand{\Omicron}{\mathrm{O}}\providecommand{\Rho}{\mathrm{P}}\providecommand{\Tau}{\mathrm{T}}
\providecommand{\Chi}{\mathrm{X}}\providecommand{\omicron}{o}\providecommand{\gt}{>}\providecommand{\lt}{<}
\providecommand{\And}{\&}\providecommand{\hdashline}{\hline}\providecommand{\operatornamewithlimits}[1]{\operatorname*{#1}}
\providecommand{\llbracket}{[\![}\providecommand{\rrbracket}{]\!]}\providecommand{\cancel}[1]{#1}
\providecommand{\mathscr}{\mathcal}
% the characters beyond ASCII this document uses
\DeclareUnicodeCharacter{2160}{\textup{I}}
\DeclareUnicodeCharacter{2161}{\textup{II}}
% the project's macros ([katex.macros] in blueprint.toml)
\providecommand{\A}{}\renewcommand{\A}{\mathbf A}
\providecommand{\AEn}{}\renewcommand{\AEn}{A_{E,n}}
\providecommand{\Aaf}{}\renewcommand{\Aaf}{\mathbf A}
\providecommand{\Ab}{}\renewcommand{\Ab}{\mathsf{Ab}}
\providecommand{\AffSch}{}\renewcommand{\AffSch}{\operatorname{AffSch}}
\providecommand{\Aone}{}\renewcommand{\Aone}{\mathbf A^1}
\providecommand{\Bl}{}\renewcommand{\Bl}{\operatorname{Bl}}
\providecommand{\Br}{}\renewcommand{\Br}{\mathrm{Br}}
\providecommand{\C}{}\renewcommand{\C}{\mathbb C}
\providecommand{\Cech}{}\renewcommand{\Cech}{\check{H}}
\providecommand{\Coext}{}\renewcommand{\Coext}{\operatorname{Coext}}
\providecommand{\Cone}{}\renewcommand{\Cone}{\operatorname{Cone}}
\providecommand{\Cons}{}\renewcommand{\Cons}{\operatorname{Cons}}
\providecommand{\Cref}{}\renewcommand{\Cref}[1]{\text{#1}}
\providecommand{\D}{}\renewcommand{\D}{D}
\providecommand{\DLAFT}{}\renewcommand{\DLAFT}{\operatorname{FT}}
\providecommand{\DLAFr}{}\renewcommand{\DLAFr}{\operatorname{Fr}}
\providecommand{\DLAQ}{}\renewcommand{\DLAQ}{{\overline{\mathbf Q}_\ell}}
\providecommand{\DLRQ}{}\renewcommand{\DLRQ}{{\overline{\mathbf Q}_\ell}}
\providecommand{\Dbc}{}\renewcommand{\Dbc}{D^b_c}
\providecommand{\Dc}{}\renewcommand{\Dc}{D_{\mathrm{cons}}}
\providecommand{\Dctf}{}\renewcommand{\Dctf}{D^b_{ctf}}
\providecommand{\Dic}{}\renewcommand{\Dic}{D_{\IndCons}}
\providecommand{\Dplus}{}\renewcommand{\Dplus}{D^{+}}
\providecommand{\End}{}\renewcommand{\End}{\operatorname{End}}
\providecommand{\Et}{}\renewcommand{\Et}{\mathrm{\acute{e}t}}
\providecommand{\Ex}{}\renewcommand{\Ex}{\mathrm{Ex}}
\providecommand{\Ext}{}\renewcommand{\Ext}{\operatorname{Ext}}
\providecommand{\F}{}\renewcommand{\F}{\mathbb F}
\providecommand{\FT}{}\renewcommand{\FT}{\operatorname{FT}}
\providecommand{\Fbar}{}\renewcommand{\Fbar}{{\overline{\mathbf F}_q}}
\providecommand{\Fell}{}\renewcommand{\Fell}{\mathbf F_\ell}
\providecommand{\Fib}{}\renewcommand{\Fib}{\operatorname{Fib}}
\providecommand{\Fix}{}\renewcommand{\Fix}{\operatorname{Fix}}
\providecommand{\Fp}{}\renewcommand{\Fp}{{\mathbf F_{p}}}
\providecommand{\Fq}{}\renewcommand{\Fq}{{\mathbf F_q}}
\providecommand{\Fqn}{}\renewcommand{\Fqn}{{\mathbf F_{q^{n}}}}
\providecommand{\Frob}{}\renewcommand{\Frob}{\operatorname{Frob}}
\providecommand{\GL}{}\renewcommand{\GL}{\operatorname{GL}}
\providecommand{\Gal}{}\renewcommand{\Gal}{\operatorname{Gal}}
\providecommand{\Hens}{}\renewcommand{\Hens}{\operatorname{Hens}}
\providecommand{\Hilb}{}\renewcommand{\Hilb}{\operatorname{Hilb}}
\providecommand{\Hom}{}\renewcommand{\Hom}{\operatorname{Hom}}
\providecommand{\IndCons}{}\renewcommand{\IndCons}{\operatorname{IndCons}}
\providecommand{\Kar}{}\renewcommand{\Kar}{\operatorname{Kar}}
\providecommand{\LT}{}\renewcommand{\LT}{\operatorname{LT}}
\providecommand{\Lie}{}\renewcommand{\Lie}{\operatorname{Lie}}
\providecommand{\Loc}{}\renewcommand{\Loc}{\operatorname{Loc}}
\providecommand{\Mod}{}\renewcommand{\Mod}{\mathsf{Mod}}
\providecommand{\N}{}\renewcommand{\N}{\mathbb{N}}
\providecommand{\Needspace}{}\renewcommand{\Needspace}{\par}
\providecommand{\O}{}\renewcommand{\O}{{\mathcal{O}}}
\providecommand{\Ob}{}\renewcommand{\Ob}{\mathrm{Ob}}
\providecommand{\Oe}{}\renewcommand{\Oe}{\mathcal O_E}
\providecommand{\PP}{}\renewcommand{\PP}{\mathbf{P}}
\providecommand{\Perf}{}\renewcommand{\Perf}{\operatorname{Perf}}
\providecommand{\Pic}{}\renewcommand{\Pic}{\mathrm{Pic}}
\providecommand{\Pm}{}\renewcommand{\Pm}{\mathbf P^{m}}
\providecommand{\Pone}{}\renewcommand{\Pone}{\mathbf P^{1}}
\providecommand{\Pqn}{}\renewcommand{\Pqn}{\mathbf P^{n}}
\providecommand{\ProFin}{}\renewcommand{\ProFin}{\operatorname{ProFin}}
\providecommand{\Proet}{}\renewcommand{\Proet}{\mathrm{pro\acute{e}t}}
\providecommand{\Q}{}\renewcommand{\Q}{\mathbb Q}
\providecommand{\Qb}{}\renewcommand{\Qb}{\overline{\mathbf Q}_{\ell}}
\providecommand{\Qbar}{}\renewcommand{\Qbar}{\overline{\mathbf Q}}
\providecommand{\QbarX}{}\renewcommand{\QbarX}{\overline{\mathbf Q}_{\ell,X}}
\providecommand{\Qbarl}{}\renewcommand{\Qbarl}{{\overline{\mathbf Q}_{\ell}}}
\providecommand{\Qell}{}\renewcommand{\Qell}{{\mathbf{Q}_{\ell}}}
\providecommand{\Ql}{}\renewcommand{\Ql}{\mathbf Q_{\ell}}
\providecommand{\Qlb}{}\renewcommand{\Qlb}{{\overline{\mathbf Q}_\ell}}
\providecommand{\Qlbar}{}\renewcommand{\Qlbar}{\overline{\mathbf Q}_\ell}
\providecommand{\R}{}\renewcommand{\R}{\mathbb{R}}
\providecommand{\RG}{}\renewcommand{\RG}{\operatorname{\mathrm{R}\Gamma}}
\providecommand{\RGamma}{}\renewcommand{\RGamma}{R\Gamma}
\providecommand{\RGammac}{}\renewcommand{\RGammac}{R\Gamma_c}
\providecommand{\RHom}{}\renewcommand{\RHom}{R\mathcal{H}om}
\providecommand{\RHomX}{}\renewcommand{\RHomX}{R\mathcal H om_X}
\providecommand{\Rat}{}\renewcommand{\Rat}{\operatorname{IsRat}}
\providecommand{\Res}{}\renewcommand{\Res}{\operatorname{Res}}
\providecommand{\Rlim}{}\renewcommand{\Rlim}{R\!\varprojlim}
\providecommand{\SIm}{}\renewcommand{\SIm}{\operatorname{SIm}}
\providecommand{\Sch}{}\renewcommand{\Sch}{\mathsf{Sch}}
\providecommand{\Sh}{}\renewcommand{\Sh}{\operatorname{Sh}}
\providecommand{\Span}{}\renewcommand{\Span}{\operatorname{span}}
\providecommand{\Spec}{}\renewcommand{\Spec}{\operatorname{Spec}}
\providecommand{\Tr}{}\renewcommand{\Tr}{\operatorname{Tr}}
\providecommand{\Trsm}{}\renewcommand{\Trsm}{\operatorname{Tr}^{\mathrm{sm}}}
\providecommand{\Wt}{}\renewcommand{\Wt}{\operatorname{Wt}}
\providecommand{\Z}{}\renewcommand{\Z}{\mathbb{Z}}
\providecommand{\Zar}{}\renewcommand{\Zar}{\mathrm{zar}}
\providecommand{\Zell}{}\renewcommand{\Zell}{{\mathbf{Z}_{\ell}}}
\providecommand{\Zl}{}\renewcommand{\Zl}{{\mathbf Z_\ell}}
\providecommand{\affet}{}\renewcommand{\affet}[1]{{#1}^{\mathrm{aff}}_{\mathrm{ét}}}
\providecommand{\affproet}{}\renewcommand{\affproet}[1]{{#1}^{\mathrm{aff}}_{\mathrm{proét}}}
\providecommand{\boxt}{}\renewcommand{\boxt}{\boxtimes_A^{\mathrm L}}
\providecommand{\cF}{}\renewcommand{\cF}{\mathcal{F}}
\providecommand{\cHom}{}\renewcommand{\cHom}{R\mathcal{H}\!om}
\providecommand{\cI}{}\renewcommand{\cI}{\mathcal{I}}
\providecommand{\cL}{}\renewcommand{\cL}{\mathcal L}
\providecommand{\cM}{}\renewcommand{\cM}{\mathcal M}
\providecommand{\cN}{}\renewcommand{\cN}{\mathcal N}
\providecommand{\cO}{}\renewcommand{\cO}{\mathcal O}
\providecommand{\cT}{}\renewcommand{\cT}{\mathcal T}
\providecommand{\checked}{}\renewcommand{\checked}{\color{red}\aftergroup\restorecolor}
\providecommand{\coeff}{}\renewcommand{\coeff}{\operatorname{coeff}}
\providecommand{\coker}{}\renewcommand{\coker}{\operatorname{coker}}
\providecommand{\colim}{}\renewcommand{\colim}{\operatorname{colim}}
\providecommand{\constructedby}{}\renewcommand{\constructedby}[1]{\par\smallskip\noindent\textcolor{purple!60!black}{\textsf{Constructed by:} }
  \footnotesize\sloppy\nolinkurl{#1}\par}
\providecommand{\deepresult}{}\renewcommand{\deepresult}[1]{\par\smallskip\noindent\textcolor{orange!45!black}{
  \textsf{Blueprint role: named deep theorem node; see the cluster proof architecture.}}
  \par\noindent\textcolor{orange!45!black}{\textsf{Mathematical source:} }#1\par}
\providecommand{\dt}{}\renewcommand{\dt}{\mathbin{\otimes^{\mathbf L}}}
\providecommand{\dw}{}\renewcommand{\dw}{\operatorname{dw}}
\providecommand{\et}{}\renewcommand{\et}[1]{{#1}_{\mathrm{ét}}}
\providecommand{\ev}{}\renewcommand{\ev}{\operatorname{ev}}
\providecommand{\fib}{}\renewcommand{\fib}{\operatorname{fib}}
\providecommand{\formalizationnote}{}\renewcommand{\formalizationnote}[1]{\begin{quote}\small\raggedright\textbf{Formalization note.} #1\end{quote}}
\providecommand{\fpqc}{}\renewcommand{\fpqc}{X_{\mathrm{fpqc}}}
\providecommand{\gammacont}{}\renewcommand{\gammacont}{\Gamma_{\mathrm{cont}}}
\providecommand{\gm}{}\renewcommand{\gm}{\mathbf G_m}
\providecommand{\hocolim}{}\renewcommand{\hocolim}{\operatorname{hocolim}}
\providecommand{\holim}{}\renewcommand{\holim}{\operatorname{holim}}
\providecommand{\hrB}{}\renewcommand{\hrB}{\@ifnextchar]{\hrEnd}{\hrStep}}
\providecommand{\hrEnd}{}\renewcommand{\hrEnd}{\@firstoftwo{}}
\providecommand{\hrStep}{}\renewcommand{\hrStep}{\@firstoftwo{\hrB}}
\providecommand{\id}{}\renewcommand{\id}{\operatorname{id}}
\providecommand{\kunFl}{}\renewcommand{\kunFl}{\mathbf F_\ell}
\providecommand{\kunZn}{}\renewcommand{\kunZn}{\mathbf Z/n\mathbf Z}
\providecommand{\kuncd}{}\renewcommand{\kuncd}{\operatorname{cd}}
\providecommand{\kundim}{}\renewcommand{\kundim}{\dim}
\providecommand{\kunotimes}{}\renewcommand{\kunotimes}[1]{\otimes^{\mathbf L}_{#1}}
\providecommand{\kunsupp}{}\renewcommand{\kunsupp}{\operatorname{Supp}}
\providecommand{\lotimes}{}\renewcommand{\lotimes}{\otimes^{\mathrm L}}
\providecommand{\m}{}\renewcommand{\m}{{\mathfrak{m}}}
\providecommand{\n}{}\renewcommand{\n}{{\mathfrak{n}}}
\providecommand{\nagataGap}{}\renewcommand{\nagataGap}[1]{\par\noindent\textbf{Formalization gap.} #1\par}
\providecommand{\openobligation}{}\renewcommand{\openobligation}[1]{\par\smallskip\noindent\textcolor{red!65!black}{\textsf{Open obligation:} }
  \footnotesize\sloppy\nolinkurl{#1}\par}
\providecommand{\p}{}\renewcommand{\p}{{\mathfrak{p}}}
\providecommand{\pr}{}\renewcommand{\pr}{\operatorname{pr}}
\providecommand{\preshv}{}\renewcommand{\preshv}[1]{\mathrm{PreShv}(#1)}
\providecommand{\pro}{}\renewcommand{\pro}[1]{\mathrm{Pro}(#1)}
\providecommand{\proet}{}\renewcommand{\proet}{X_{\mathrm{proét}}}
\providecommand{\q}{}\renewcommand{\q}{{\mathfrak{q}}}
\providecommand{\rad}{}\renewcommand{\rad}{\operatorname{rad}}
\providecommand{\rank}{}\renewcommand{\rank}{\operatorname{rank}}
\providecommand{\red}{}\renewcommand{\red}{\color{red}}
\providecommand{\restorecolor}{}\renewcommand{\restorecolor}{\color{black}}
\providecommand{\retiredobligation}{}\renewcommand{\retiredobligation}[1]{\par\smallskip\noindent\textcolor{red!75!black}{\textsf{Retired invalid route:} }
  \footnotesize\sloppy\nolinkurl{#1}\par}
\providecommand{\rk}{}\renewcommand{\rk}{\operatorname{rk}}
\providecommand{\sep}{}\renewcommand{\sep}{\textnormal{sep}}
\providecommand{\shv}{}\renewcommand{\shv}[1]{\mathrm{Shv}(#1)}
\providecommand{\spec}{}\renewcommand{\spec}[1]{\mathrm{Spec}(#1)}
\providecommand{\stacks}{}\renewcommand{\stacks}[1]{\href{https://stacks.math.columbia.edu/tag/#1}{\textnormal{Stacks, Tag~#1}}}
\providecommand{\stacksproject}{}\renewcommand{\stacksproject}[1]{{\cite[\href{https://stacks.math.columbia.edu/tag/#1}{Tag #1}]{stacks-project}}}
\providecommand{\statusnote}{}\renewcommand{\statusnote}[1]{\begin{quote}\small\raggedright\textbf{Status.} #1\end{quote}}
\providecommand{\tildeideal}{}\renewcommand{\tildeideal}[2]{\widetilde{(#1, #2)}}
\providecommand{\tildering}{}\renewcommand{\tildering}[3]{{#1}_{\widetilde{(#2, #3)}}}
\providecommand{\tr}{}\renewcommand{\tr}{\operatorname{tr}}
\providecommand{\wt}{}\renewcommand{\wt}{\operatorname{wt}}

\begin{document}
\pagestyle{plain}
\begin{center}
{\LARGE\bfseries Weil Conjectures\par}
\vspace{0.7em}
{\large Chapter 6: Constructibility of $f^!$\par}
\vspace{0.5em}
{\small\color{bpmuted}Snapshot 2026-09-26, revision 8e4cb88e2a\par}
\end{center}
\vspace{0.5em}
{\small\tableofcontents}
\bigskip\hrule\bigskip

% object: chap-Constructibility-of-f (blueprint/chap-Constructibility-of-f/_section.md)
\section*{\texorpdfstring{6\quad Constructibility of $f^!$}{6 Constructibility of f!}}
\phantomsection\label{bp1}%
\addcontentsline{toc}{section}{\texorpdfstring{6\quad Constructibility of $f^!$}{6 Constructibility of f!}}
% text: prose-06-01
\phantomsection\label{bp2}%
Brown representability gave us \(f^!\) (Theorem \bpout{5.31}), but only as a functor

\[
f^!:\Loc(D^b_{\mathrm{ctf}}(Y,\Lambda))\longrightarrow\Loc(D^b_{\mathrm{ctf}}(X,\Lambda))
\]

for \(f:X\to Y\) separated of finite type and \(\Lambda\) an admissible coefficient ring. It tells us nothing about the values of \(f^!\). A priori, \(f^!\) of a constructible complex is just some object of the large category on the right, possibly unbounded and with cohomology sheaves that are not constructible.

The rest of the notes needs more than this. For example, Verdier duality (Chapter \bpout{10}) says that \(\mathbb D_X=\cHom(-,a_X^!\Lambda)\) is an involution of \(D^b_{\mathrm{ctf}}(X,\Lambda)\), and this makes sense only if the dualizing complex \(a_X^!\Lambda\) lies in \(D^b_{\mathrm{ctf}}(X,\Lambda)\). Similarly, the local terms of the trace formula are defined through the isomorphism \(\cHom(c_2^*K,c_1^!K)\simeq c^!(K\boxtimes\mathbb D_XK)\) for a correspondence \(c=(c_1,c_2)\) (Lemma \bpout{11.7}), which also needs \(f^!\) to preserve constructible complexes. In this chapter we prove what both need (part (4) of Theorem \bpout{5.2}): \(f^!\) maps \(D^b_{\mathrm{ctf}}(Y,\Lambda)\) into \(D^b_{\mathrm{ctf}}(X,\Lambda)\). A universal property cannot give this, so we have to compute \(f^!\), and the computation is local.

By Proposition \bpout{5.32}, the functor \(f^!\) is compatible with composition, equals \(j^*\) for an open immersion \(j\), and is local on \(X\) and on \(Y\). We may therefore assume that \(Y\) is affine and that \(f\) is a closed immersion into some \(\mathbf A^n_Y\) followed by the projection to \(Y\). Since \(\mathbf A^n_Y\) is open in \((\mathbf P^1_Y)^n\), and \((\mathbf P^1_Y)^n\to Y\) is a composite of \(n\) projections \(\mathbf P^1_T\to T\), two cases remain.

For the projection \(p:\mathbf P^1_T\to T\) we can write the adjunction down, and we find \(p^!=p^*(1)[2]\) (Proposition \hyperref[bp22]{6.7}). Here the counit is the trace \(Rp_*\Lambda(1)[2]\to\Lambda\), and the unit is the cycle class of the diagonal in \(\mathbf P^1_T\times_T\mathbf P^1_T\). The triangle identities then hold because the diagonal meets each fibre of either projection in a single point. For a closed immersion \(i:Z\hookrightarrow X\) with open complement \(j:U\hookrightarrow X\), on the other hand, the functor \(i^!\) takes sections with support in \(Z\), and it sits in the localization triangle

\[
i_*i^!K\longrightarrow K\longrightarrow Rj_*j^*K\longrightarrow .
\]

Let \(K\) be constructible. Then \(j^*K\) is constructible, and since \(i^*i_*\simeq\mathrm{id}\) and \(i_*\) preserves constructibility, \(i^!K\) is constructible exactly when the third term \(Rj_*j^*K\) is. So the constructibility of \(f^!\) comes down to one question: does direct image along an open immersion preserve bounded constructible complexes and, with finite coefficients, finite Tor amplitude? Proper base change does not help here, since \(j\) is not proper.

The answer is not formal, because \(Rj_*\), and with it \(i^!\), sees more than the stalks of \(K\) on \(Z\): it sees the cohomology of \(K\) on a punctured neighbourhood of \(Z\). For the origin \(i:0\hookrightarrow\mathbf A^1\) over an algebraically closed field, the stalk of \(Rj_*\Lambda\) at \(0\) is the cohomology of a punctured disc, which the Kummer covers of \(\mathbf G_m\) compute to be \(\Lambda\oplus\Lambda(-1)[-1]\), and the triangle gives \(i^!\Lambda\simeq\Lambda(-1)[-2]\), whereas \(i^*\Lambda=\Lambda\) (Example \hyperref[bp79]{6.25} below).

Over a field the answer is Deligne's finiteness theorem in SGA 4\(\tfrac12\) [\hyperref[bp89]{SGA4½}], and we follow his proof. Its induction on dimension leads to fibrations over bases of positive dimension, so Deligne proves a statement in families over a base \(S\). In families, however, the formation of \(Rj_*\) need not commute with base change. Let \(S=\mathbf A^1\) with coordinate \(t\), let \(Y=\mathbf A^2\) with coordinates \((t,x)\), viewed over \(S\), and let \(j\) be the inclusion of \(X=Y\setminus\{(0,0)\}\). By purity (Theorem \bpout{9.3}, proved in Chapter \bpout{9} below), the punctured strict henselization of \(Y\) at the origin has cohomology \(\Lambda\) in degree \(0\) and \(\Lambda(-2)\) in degree \(3\) only. Hence \(R^1j_*\Lambda=0\). On the fibre over \(t=0\), however, \(j\) becomes \(\mathbf G_m\hookrightarrow\mathbf A^1\), whose \(R^1j_*\Lambda\) is \(\Lambda(-1)\) at the origin. Since over \(t\neq0\) the map \(j\) is an isomorphism, base change fails only at \(t=0\).

Deligne's generic base change theorem says that this is the general picture. Let \(j:X\hookrightarrow Y\) be an open immersion of schemes of finite type over an integral base \(S\) of finite type over \(k\), and let \(F\) be a constructible \(\Fell\)-sheaf on \(X\). Then there is a nonempty open \(S_0\subset S\) over which \(Rj_*F\) is bounded constructible and its formation commutes with base change along every \(T\to S_0\) of finite type (Theorem \hyperref[bp51]{6.17}; Deligne proves more). Over \(S=\operatorname{Spec}k\) the only nonempty open is \(S\) itself, so we obtain the finiteness and base change of \(Rj_*\) over a field (Corollary \hyperref[bp54]{6.18}).

The projective line returns in the proof of the generic theorem. There we must control \(Rf_*L\) for smooth \(f\) and locally constant \(L\), and we do so through the formula for \(f^!\) of a smooth morphism (Section \hyperref[bp24]{6.3}).

For finite coefficients (Theorem \hyperref[bp61]{6.20}) we also bound the Tor amplitude of \(Rj_*K\): it grows by at most \(\operatorname{cd}_\ell(j)\) (Proposition \hyperref[bp58]{6.19}). For integral and rational coefficients we have \(D^b_{\mathrm{ctf}}=\Dbc\), and we reduce modulo a uniformizer. This works because an integral complex is constructible if it is derived complete with constructible reduction (Proposition \bpout{3.11}). Only the adjunction for the projective line resists this reduction. Its unit and counit must be constructed as single morphisms with \(\mathcal O_E\)-coefficients (Proposition \hyperref[bp69]{6.22}). The result is Theorem \hyperref[bp84]{6.27}.

% object: sec-Conventions-2 (blueprint/chap-Constructibility-of-f/sec-Conventions-2/_section.md), content hidden
\subsection*{\texorpdfstring{6.1\quad Conventions}{6.1 Conventions}}
\phantomsection\label{bp3}%
\addcontentsline{toc}{subsection}{\texorpdfstring{6.1\quad Conventions}{6.1 Conventions}}
% text: prose-06-02
\phantomsection\label{bp4}%
Until Section \hyperref[bp56]{6.6} the coefficient ring is a finite commutative \(\ell\)-primary ring, which we denote by \(A\). With such coefficients nothing is lost by computing on the étale site. For \(X\) of finite type over \(k\), the comparison functor

\[
\nu_X^*:D(X_{\acute{e}t},A)\longrightarrow D(X_{\mathrm{pro}\acute{e}t},A)
\]

is fully faithful (Proposition \bpout{2.49}), and it identifies \(D(X_{\acute{e}t},A)\) with the localizing category \(\Loc(D^b_{\mathrm{ctf}}(X,A))\) on which \(f^!\) was constructed (Proposition \bpout{5.15}). It also identifies the constructible complexes on the two sites (Theorem \bpout{2.43}). In the following sections we therefore work in \(D(X_{\acute{e}t},A)\), and every result is transported to the pro-étale site by \(\nu_X^*\).

This transport is harmless because \(\nu^*\) identifies all the operations we use with their étale counterparts, compatibly with units and counits (Lemma \bpout{2.39}). These operations are the cohomology (Theorem \bpout{2.35}); the direct images and their base change maps, also for unbounded complexes (Theorem \bpout{2.38} and Proposition \bpout{2.49}); and the inverse images, tensor products, extension by zero and the localization triangles (Section \bpout{2.9}). Moreover, by Corollary \bpout{2.44}, finite Tor amplitude can be checked on the étale site.

We write \(\Dbc(X,A)\) also for the bounded constructible objects of \(D(X_{\acute{e}t},A)\), which correspond to those of the pro-étale category. For a separated morphism of finite type \(f:X\to S\), the functor \(Rf_!\) and its right adjoint \(f^!\) of Theorem \bpout{5.31} correspond to the functors \(Rf_!\) and \(f^!\) between \(D(X_{\acute{e}t},A)\) and \(D(S_{\acute{e}t},A)\), and the two adjunctions correspond. Accordingly, from now on \(g^*\), \(Rf_*\), \(Rf_!\), \(f^!\) and \(\otimes^{\mathbf L}\) denote the functors on the étale derived categories.

Below we apply Proposition \bpout{5.28}, stated on \(\Loc(D^b_{\mathrm{ctf}}(-,A))\), to the étale functors. This is legitimate because the comparison identifies the composition maps of Proposition \bpout{5.28} on the two sites: the refinement maps (5.7) and the open–proper interchange (5.9) are characterized by their restriction to the open. The comparison also preserves the counits defining base change and the inverse counits defining forgetting supports (Proposition \bpout{5.29}). Taking right adjoints, we see that it identifies the composition maps for \(f^!\) as well.

Except in Lemmas \hyperref[bp11]{6.3}, \hyperref[bp14]{Lemma 6.4} and \hyperref[bp41]{Lemma 6.13}, which allow more general bases as stated there, all schemes and base changes are of finite type over \(k\).

We use three classical theorems about the étale site.

% object: theorem-results-from-etale-theory (blueprint/chap-Constructibility-of-f/sec-Conventions-2/theorem-results-from-etale-theory/_section.md)
\phantomsection\label{bp5}%
\begin{bpstatementit}{Theorem~6.1}{ (Results from the étale theory)}
Let \(f:X\to S\) be a morphism of schemes of finite type over \(k\).

\begin{enumerate}
\item (Proper base change) If \(f\) is proper, then for every morphism of schemes \(g:S'\to S\), not necessarily of finite type over \(k\), and every \(K\in D(X_{\acute{e}t},A)\), the base change map \(g^*Rf_*K\to Rf'_*g'^*K\) is an isomorphism, where \(f':X'\to S'\) and \(g':X'\to X\) are the projections of \(X'=X\times_SS'\). In particular, the geometric stalks of \(Rf_*K\) are the cohomology of the geometric fibres.
\item (Finiteness) If \(f\) is separated, then \(Rf_!\) maps \(\Dbc(X,A)\) into \(\Dbc(S,A)\).
\item (Cohomological dimension) The functors \(Rf_*\) and, if \(f\) is separated, \(Rf_!\) have finite cohomological amplitude on torsion sheaves. In particular \(\operatorname{cd}_\ell(f)<\infty\) in the sense of Definition \bpout{2.45}.
\end{enumerate}
\end{bpstatementit}

% proof: proof/theorem-results-from-etale-theory
\begin{proof}[{Proof}]\phantomsection\label{bp6}%
For bounded below complexes, (1) is Theorem \bpout{4.33} together with the comparison of direct images (Theorem \bpout{2.38}); for unbounded complexes it is [Stacks, \href{https://stacks.math.columbia.edu/tag/0F0C}{0F0C}], whose finiteness hypothesis on \(f\) is part (3). Part (2) is [Stacks, \href{https://stacks.math.columbia.edu/tag/0GL0}{0GL0}]. As for part (3), for \(Rf_*\) it is [Stacks, \href{https://stacks.math.columbia.edu/tag/0F10}{0F10}], and for \(Rf_!\) we write \(f=\bar f\circ j\) with \(\bar f\) proper and \(j\) an open immersion; then \(Rf_!=R\bar f_*j_!\) with \(j_!\) exact, and [Stacks, \href{https://stacks.math.columbia.edu/tag/0F10}{0F10}] applies to \(\bar f\). Finally, the last assertion follows by the comparison of direct images, since classical sheaves are of the form \(\nu_X^*F\).
\end{proof}

% object: definition-exceptional-base-change-finite (blueprint/chap-Constructibility-of-f/sec-Conventions-2/definition-exceptional-base-change-finite/_section.md)
\phantomsection\label{bp7}%
\begin{bpstatement}{Definition~6.2}{ (The base change map for $f^!$, finite coefficients)}
Let \(f:X\to S\) be separated and of finite type, let \(g:S'\to S\) be a morphism, let \(X'=X\times_SS'\), and let \(f':X'\to S'\), \(g':X'\to X\) be the projections. For \(N\in D(S_{\acute{e}t},A)\), let \(\varepsilon_N:Rf_!f^!N\to N\) be the counit of \(Rf_!\dashv f^!\). The \emph{base change map}

\[
\beta_N:g'^*f^!N\longrightarrow f'^!g^*N
\tag{6.1}
\]

is the transpose, for \(f'\), of the morphism

\[
Rf'_!g'^*f^!N
\xrightarrow{\ \sim\ }g^*Rf_!f^!N
\xrightarrow{g^*\varepsilon_N}g^*N,
\tag{6.2}
\]

where the first arrow is the inverse of the base change isomorphism for \(Rf_!\) (Proposition proposition-composition-base-change-rf-shriek).
\end{bpstatement}

% text: prose-06-03
\phantomsection\label{bp8}%
It is an isomorphism for smooth \(f\) (Proposition \hyperref[bp30]{6.9}). The global isomorphism \(f^!\simeq f^*(d)[2d]\) for smooth \(f\) of relative dimension \(d\), and its compatibility with composition, are treated with smooth purity in Chapter \bpout{9}. Here we need only its restriction to charts \(U\to\mathbf A^d_S\).


% object: sec-Duality-for-the-projective-line (blueprint/chap-Constructibility-of-f/sec-Duality-for-the-projective-line/_section.md), content hidden
\subsection*{\texorpdfstring{6.2\quad Duality for the projective line}{6.2 Duality for the projective line}}
\phantomsection\label{bp9}%
\addcontentsline{toc}{subsection}{\texorpdfstring{6.2\quad Duality for the projective line}{6.2 Duality for the projective line}}
% text: prose-06-04
\phantomsection\label{bp10}%
Let \(p:\mathbf P^1_S\to S\) be the projection. We construct the adjunction \(Rp_*\dashv p^*(1)[2]\) from two maps: the counit is a trace \(Rp_*A(1)[2]\to A\), obtained from the cohomology of \(\mathbf P^1\), and the unit is the cycle class of the diagonal in \(\mathbf P^1_S\times_S\mathbf P^1_S\), a special case of the cycle class of an effective Cartier divisor. Both triangle identities say that pushing the class of the diagonal forward along either projection \(\mathbf P^1\times\mathbf P^1\to\mathbf P^1\) and applying the trace gives \(1\). This holds because the diagonal meets each fibre of either projection in a single point.

In topology the same two maps appear. For instance, for a bundle \(E\to B\) with fibre \(S^2\), integration over the fibre \(H^2(S^2)\to\mathbf Z\) is the counit of Poincaré duality on the fibres, and the unit is the Thom class of the diagonal of \(E\times_BE\), which is Poincaré dual to the identity correspondence. In étale cohomology, by contrast, the fundamental class of a point on \(\mathbf P^1\) comes from a line bundle through the Kummer sequence. So both maps are built from line bundles: the trace comes from \(c_1(\mathcal O(1))\), and the class of the diagonal comes from the divisor \(x_0y_1-x_1y_0=0\).

For every other morphism, we shall compute \(f^!\) from this adjunction, from \(j_!\dashv j^*\) for open immersions and from localization.

% object: lemma-cycle-class-cartier-divisor (blueprint/chap-Constructibility-of-f/sec-Duality-for-the-projective-line/lemma-cycle-class-cartier-divisor/_section.md)
\phantomsection\label{bp11}%
\begin{bpstatementit}{Lemma~6.3}{ (The cycle class of a Cartier divisor)}
Let \(\Lambda\) be a finite commutative \(\ell\)-primary ring, and let \(i:D\hookrightarrow Z\) be an effective Cartier divisor on a scheme \(Z\) on which \(\ell\) is invertible. There is a morphism

\[
c_D:i_*\Lambda_D\longrightarrow\Lambda_Z(1)[2]\quad\text{in }D(Z_{\acute{e}t},\Lambda),
\]

and hence, via \(\nu_Z^*\), in \(D(Z_{\mathrm{pro}\acute{e}t},\Lambda)\), whose composite with \(\Lambda_Z\to i_*\Lambda_D\) is the Kummer first Chern class \(c_1(\mathcal O_Z(D))\). It is compatible with change of coefficients and with every base change along which the pullback of the divisor is still an effective Cartier divisor.
\end{bpstatementit}

% proof: proof/lemma-cycle-class-cartier-divisor
\begin{proof}[{Proof}]\phantomsection\label{bp12}%
The map \(c_D\) will be the boundary map of the extension of \(i_*\mathbf Z_D\) by \(\mathbf G_m\) formed by meromorphic functions with poles along \(D\), reduced modulo \(\ell^r\); this extension is the étale version of the divisor sequence \(0\to\mathcal O_Z\to\mathcal O_Z(D)\to
\mathcal O_Z(D)|_D\to0\).

We construct \(c_D\) in \(D(Z_{\acute{e}t},\Lambda)\). The functor \(\nu_Z^*\), which commutes with inverse images and with change of coefficients, then transports it and its compatibilities to \(Z_{\mathrm{pro}\acute{e}t}\).

Let \(j:Z\setminus D\hookrightarrow Z\) be the open complement. Let \(\mathcal M_D\subset j_*\mathbf G_m\) be the subsheaf which, on an open where \(D\) has a regular equation \(t\), consists of the products \(ut^{-m}\) with \(u\) a unit and \(m\) a locally constant integer on \(D\); equivalently, it is generated by \(\mathbf G_m\) and \(t^{-1}\). Since replacing \(t\) by a unit multiple gives the same subsheaf, \(\mathcal M_D\) is globally defined. Sending \(t^{-1}\) to \(1\), we obtain a short exact sequence of abelian sheaves on \(Z_{\acute{e}t}\)

\[
\Sigma_D:\ 0\longrightarrow\mathbf G_m\longrightarrow \mathcal M_D
\longrightarrow i_*\mathbf Z_D\longrightarrow0,
\tag{6.3}
\]

whose exactness can be checked on geometric stalks. Away from \(D\), the function \(t\) is a unit, so there \(\mathcal M_D=\mathbf G_m\), and both the quotient and the stalk of \(i_*\mathbf Z_D\) are zero. At a point of \(D\), on the other hand, no nonzero power of a regular nonunit \(t\) is a unit, and therefore the exponent \(m\) in \(ut^{-m}\) is well defined. The kernel of the map \(ut^{-m}\mapsto m\) is then \(\mathbf G_m\), and the map is surjective because exponents lift locally to powers of \(t^{-1}\). For example, for \(D=\{0\}\subset Z=\mathbf A^1=\operatorname{Spec}k[t]\), the sections of \(\mathcal M_D\) near the origin are the functions \(ut^{-m}\), that is, the rational functions whose only zero or pole is at \(0\), and the map records the order of the pole, \(m=-\operatorname{ord}_0\).

The boundary of \(\Sigma_D\) is a morphism \(i_*\mathbf Z_D\to\mathbf G_m[1]\) in \(D(Z_{\acute{e}t},\mathbf Z)\), to which we apply \(-\otimes^{\mathbf L}_{\mathbf Z}\mathbf Z/\ell^r\mathbf Z\). The source becomes \(i_*\mathbf Z/\ell^r\mathbf Z\). For the target, the Kummer sequence [Stacks, \href{https://stacks.math.columbia.edu/tag/03PL}{03PL}] shows that the cone of multiplication by \(\ell^r\) on \(\mathbf G_m\) has the single cohomology sheaf \(\mu_{\ell^r}\) in degree \(-1\), and hence the target becomes \(\mu_{\ell^r}[2]\). This is \(c_D\) for \(\mathbf Z/\ell^r\mathbf Z\). For a finite \(\mathbf Z/\ell^r\mathbf Z\)-algebra \(\Lambda\), we apply \(-\otimes^{\mathbf L}_{\mathbf Z/\ell^r}\Lambda\). Since \(\mu_{\ell^r}\) is locally free of rank one over \(\mathbf Z/\ell^r\mathbf Z\), the target becomes \(\Lambda_Z(1)[2]\).

To identify the composite with \(c_1\), we cover \(Z\) by opens with local equations \(t_\alpha\). The local lifts of the section \(1\) of \(i_*\mathbf Z_D\) through \(\Sigma_D\) are the \(t_\alpha^{-1}\). Their ratios are the transition functions of \(\mathcal O_Z(D)\); this fixes the sign convention. Hence the boundary of \(\Sigma_D\) followed by the Kummer boundary is the class of \(\mathcal O_Z(D)\). Changing a local equation multiplies a lift by a unit, and this changes the cocycle by a coboundary.

Naturality holds because the local equations, the units and the sequence \(\Sigma_D\) pull back, compatibly with the boundary maps. On the source we use the identity \(g^*i_*=i'_*g_D^*\) for a base change \(g:Z'\to Z\) with \(i':D'\hookrightarrow Z'\) and \(g_D:D'\to D\); it can be checked on geometric stalks.
\end{proof}

% text: prose-06-05
\phantomsection\label{bp13}%
The lemma will be applied twice: to the section at infinity of \(\mathbf P^1_S\), which gives the trace, and to the diagonal of \(\mathbf P^1_S\times_S\mathbf P^1_S\), which gives the unit. The first application enters the computation of \(Rp_*A\).

% object: lemma-cohomology-projective-line-trace (blueprint/chap-Constructibility-of-f/sec-Duality-for-the-projective-line/lemma-cohomology-projective-line-trace/_section.md)
\phantomsection\label{bp14}%
\begin{bpstatementit}{Lemma~6.4}{ (Cohomology of the projective line and the trace)}
Let \(\Lambda\) be a finite commutative \(\ell\)-primary ring, let \(S\) be a Noetherian scheme on which \(\ell\) is invertible, and let \(p:\mathbf P^1_S\to S\) be the projection. The unit \(\Lambda_S\to Rp_*\Lambda_{\mathbf P^1_S}\) and the class \(c_1(\mathcal O_{\mathbf P^1_S}(1))\) define an isomorphism

\[
\sigma_S:\Lambda_S\oplus\Lambda_S(-1)[-2]
\ \xrightarrow{\ \sim\ }\ Rp_*\Lambda_{\mathbf P^1_S}
\quad\text{in }D^+(S_{\acute{e}t},\Lambda),
\tag{6.4}
\]

compatible with base change in the sense that for \(g:S'\to S\) with \(S'\) Noetherian, the base change map \(g^*Rp_*\Lambda\to Rp'_*\Lambda\) carries \(g^*\sigma_S\) to \(\sigma_{S'}\). Let

\[
\operatorname{Tr}_p:Rp_*\Lambda_{\mathbf P^1_S}(1)[2]\longrightarrow\Lambda_S
\tag{6.5}
\]

be the morphism given by the projection onto the summand \(\Lambda_S\) of the twist of \(\sigma_S\). Then \(\operatorname{Tr}_p\) is compatible with every Noetherian base change.
\end{bpstatementit}

% proof: proof/lemma-cohomology-projective-line-trace
\begin{proof}[{Proof}]\phantomsection\label{bp15}%
By the comparison theorems for the cohomology and the proper direct image of bounded below complexes (Theorems \bpout{2.35} and \bpout{Theorem 2.38}), it suffices to prove the assertions on \(\sigma_S\) and \(\operatorname{Tr}_p\) in \(D^+(S_{\acute{e}t},\Lambda)\) for Noetherian \(S\).

Both components of \(\sigma_S\) are globally defined. The second comes from Lemma \hyperref[bp11]{6.3} applied to the section at infinity, whose line bundle is \(\mathcal O(1)\). Geometric stalks detect isomorphisms, and by proper base change the stalk of \(Rp_*\Lambda\) at a geometric point \(\bar s\) is \(R\Gamma(\mathbf P^1_{\bar s},\Lambda)\). It therefore suffices to check that \(\sigma_S\) is an isomorphism on a geometric fibre, first for \(\Lambda=\mathbf Z/\ell^r\mathbf Z\). By the genus zero case of Proposition \bpout{D.3}, the cohomology lives only in degrees \(0\) and \(2\), with \(H^0=\mathbf Z/\ell^r\mathbf Z\) and \(H^2=(\mathbf Z/\ell^r\mathbf Z)(-1)\). The Kummer description sends a line bundle to its degree modulo \(\ell^r\). As \(\mathcal O(1)\) has degree one, its class generates \(H^2\), and \(\sigma_S\) is an isomorphism. The unit and the class of \(\mathcal O(1)\) commute with base change. Therefore so do \(\sigma_S\), its inverse, and \(\operatorname{Tr}_p\).

For a finite \(\mathbf Z/\ell^r\mathbf Z\)-algebra \(\Lambda\), the projection formula [Stacks, \href{https://stacks.math.columbia.edu/tag/0GL5}{0GL5}] gives \(Rp_*\Lambda=Rp_*(\mathbf Z/\ell^r\mathbf Z)
\otimes^{\mathbf L}_{\mathbf Z/\ell^r}\Lambda\). Under this identification, \(\sigma_S\) for \(\Lambda\) is \(\sigma_S\) for \(\mathbf Z/\ell^r\mathbf Z\) tensored with \(\Lambda\). It is therefore again an isomorphism compatible with base change.
\end{proof}

% text: prose-06-06
\phantomsection\label{bp16}%
Over \(S=\operatorname{Spec}\bar k\), (6.4) reads \(R\Gamma(\mathbf P^1_{\bar k},A)\simeq A\oplus A(-1)[-2]\). Here the shift \([-2]\) places the second summand in degree \(2\), so that it is \(H^2\), and the twist \((-1)\) records how Galois acts on \(H^2\): the Kummer map identifies \(H^2\) with \(\operatorname{Pic}(\mathbf P^1)\otimes\mu_{\ell^r}^{\otimes-1}\), that is, with \(A(-1)\). Twisting by \((1)\) and shifting by \([2]\) therefore moves this summand to degree \(0\) with trivial twist. The map \(\operatorname{Tr}_p\) is the projection onto this summand, and it sends the class of a point to \(1\). It is the counit of \(Rp_*\dashv p^*(1)[2]\) in Proposition \hyperref[bp22]{6.7}. After its lift to \(\mathcal O_E\) in Subsection \hyperref[bp67]{6.7.1}, it is also the trace map \(\operatorname{Tr}_p\) of \(\mathbf P^1\) in Section \bpout{9.6}.

As a check of the twist, let us work over \(\mathbf F_q\). The geometric Frobenius acts by \(1\) on \(H^0\) and by \(q\) on \(H^2=A(-1)\). The alternating sum of the traces is \(1+q=\#\mathbf P^1(\mathbf F_q)\), as in Corollary \bpout{12.31}.

% object: example-compact-supports-running-example (blueprint/chap-Constructibility-of-f/sec-Duality-for-the-projective-line/example-compact-supports-running-example/_section.md)
\phantomsection\label{bp17}%
\begin{bpstatement}{Example~6.5}{ (Compact supports on the running example)}
We use the open immersion \(j_\infty:\mathbf A^1\to\mathbf P^1\) of Example \bpout{3.19}, base changed to \(\bar k\), and we let \(i_\infty:\{\infty\}\to\mathbf P^1_{\bar k}\) be the inclusion of the complementary point.

By definition, \(R\Gamma_c(\mathbf A^1,A)=R\Gamma(\mathbf P^1,j_{\infty!}A)\). The localization triangle \(j_{\infty!}A\to A_{\mathbf P^1}\to i_{\infty*}A\) and (6.4) show that

\[
R\Gamma_c(\mathbf A^1_{\bar k},A)
\simeq\fib\bigl(A\oplus A(-1)[-2]\to A\bigr)
\simeq A(-1)[-2],
\]

since restriction to \(\infty\) is an isomorphism on \(H^0\). In other words, \(H^2_c(\mathbf A^1)=A(-1)\) is the only nonzero compactly supported cohomology group. The trace on it is the same \(\operatorname{Tr}_p\).
\end{bpstatement}

% text: prose-06-07
\phantomsection\label{bp18}%
\textbf{Exercise.} Using the localization sequence of Example \bpout{3.19}, show that \(H^1_c(\mathbf G_{m,\bar k},A)\simeq A\) and \(H^2_c(\mathbf G_{m,\bar k},A)\simeq A(-1)\), and that these are the only nonzero groups.

The adjunction is built from the trace \(\operatorname{Tr}_p\) of (6.5) and the cycle class of the diagonal provided by Lemma \hyperref[bp11]{6.3}. As in all of this part of the section, we compute on the étale sites (Section \hyperref[bp3]{6.1}). Let \(p:\mathbf P^1_S\to S\) be the projection, let \(p_1,p_2:\mathbf P^1_S\times_S\mathbf P^1_S\to\mathbf P^1_S\) be the two projections and let \(\Delta:\mathbf P^1_S\hookrightarrow\mathbf P^1_S\times_S\mathbf P^1_S\) be the diagonal. For \(i=1,2\), the projection \(p_i\) is a projective line over the other factor, with trace \(\operatorname{Tr}_{p_i}\) given by (6.5). Both triangle identities rest on the following computation.

% object: lemma-class-diagonal-projective-line (blueprint/chap-Constructibility-of-f/sec-Duality-for-the-projective-line/lemma-class-diagonal-projective-line/_section.md)
\phantomsection\label{bp19}%
\begin{bpstatementit}{Lemma~6.6}{ (The class of the diagonal of $\mathbf P^1_S$)}
Let \(S\) be of finite type over \(k\), and let \(p_i\) and \(\operatorname{Tr}_{p_i}\) be as above. The diagonal \(\Delta\) is a relative effective Cartier divisor in \(\mathbf P^1_S\times_S\mathbf P^1_S\) over \(S\), and it remains one after every base change \(S'\to S\). Let

\[
\delta:\Delta_*A_{\mathbf P^1_S}\longrightarrow
A_{\mathbf P^1_S\times_S\mathbf P^1_S}(1)[2]
\]

be its cycle class (Lemma \hyperref[bp11]{6.3}). Then for \(i=1,2\) the composite

\[
s_i:A_{\mathbf P^1_S}=Rp_{i*}\Delta_*A_{\mathbf P^1_S}
\xrightarrow{Rp_{i*}\delta}
Rp_{i*}A_{\mathbf P^1_S\times_S\mathbf P^1_S}(1)[2]
\xrightarrow{\operatorname{Tr}_{p_i}}A_{\mathbf P^1_S}
\tag{6.6}
\]

is the identity.
\end{bpstatementit}

% proof: proof/lemma-class-diagonal-projective-line
\begin{proof}[{Proof}]\phantomsection\label{bp20}%
Let \(x_0,x_1\) and \(y_0,y_1\) be the pullbacks along \(p_1\) and \(p_2\) of the homogeneous coordinates of \(\mathbf P^1_S\). Then \(\Delta\) is the zero locus of the section \(x_0y_1-x_1y_0\) of \(p_1^*\mathcal O(1)\otimes p_2^*\mathcal O(1)\). On the open where \(x_0\) and \(y_0\) are invertible, in the affine coordinates \(x=x_1/x_0\) and \(y=y_1/y_0\), this equation becomes \(y-x\). It is a nonzerodivisor, and the quotient \(\mathcal O_S[x,y]/(y-x)\simeq\mathcal O_S[x]\) is flat over \(S\). The open where \(x_1\) and \(y_1\) are invertible is treated in the same way, and these two opens cover \(\Delta\). Hence \(\Delta\) is a relative effective Cartier divisor, and it remains one after every base change.

To show \(s_i=1\), note that both sides are endomorphisms of a sheaf placed in degree zero. We have

\[
\Hom_{D((\mathbf P^1_S)_{\acute{e}t},A)}
(A_{\mathbf P^1_S},A_{\mathbf P^1_S})
=\Hom_{\operatorname{Mod}((\mathbf P^1_S)_{\acute{e}t},A)}
(A_{\mathbf P^1_S},A_{\mathbf P^1_S}),
\]

so \(s_i\) is a morphism of sheaves and is determined by its geometric stalks. By proper base change and the compatibility of \(\delta\) and \(\operatorname{Tr}_{p_i}\) with base change, the stalk of (6.6) at a geometric point \(\bar x\) of \(\mathbf P^1_S\) is computed on the fibre \(p_i^{-1}(\bar x)\simeq\mathbf P^1_{\bar x}\). On this fibre the equation \(x_0y_1-x_1y_0\) restricts to a nonzero linear form, so the diagonal cuts out a single reduced point, whose line bundle is \(\mathcal O(1)\). By Lemma \hyperref[bp11]{6.3}, the class of this point composed with the unit is \(c_1(\mathcal O(1))\). In turn, the trace \(\operatorname{Tr}_p\) sends \(c_1(\mathcal O(1))\) to \(1\), since by definition (6.5) it is the projection onto the summand of (6.4) spanned by this class. Thus the stalk of \(s_i\) at every geometric point \(\bar x\) sends \(1\) to \(1\), and so \(s_i=1\).
\end{proof}

% text: prose-06-08
\phantomsection\label{bp21}%
We build the unit and counit of the adjunction from \(\delta\) and from the trace. For \(M\in D((\mathbf P^1_S)_{\acute{e}t},A)\), let

\[
\begin{split}
\eta_M:M
&=Rp_{2*}\Delta_*M\\
&\longrightarrow Rp_{2*}\bigl(p_1^*M(1)[2]\bigr)
\xrightarrow{\ \sim\ }p^*Rp_*M(1)[2] .
\end{split}
\tag{6.7}
\]

Here the first arrow is \(Rp_{2*}\) applied to \(\delta\otimes p_1^*M\), where we use the projection formula \(\Delta_*M=\Delta_*A\otimes^{\mathbf L}p_1^*M\), and the second arrow is the proper base change isomorphism \(p^*Rp_*\simeq Rp_{2*}p_1^*\). For \(N\in D(S_{\acute{e}t},A)\), let

\[
\varepsilon_N:Rp_*\bigl(p^*N(1)[2]\bigr)
\xrightarrow{\ \sim\ }N\otimes_A^{\mathbf L}
Rp_*A_{\mathbf P^1_S}(1)[2]
\xrightarrow{1\otimes \operatorname{Tr}_p}N,
\tag{6.8}
\]

where the first arrow is the projection formula. These maps exist for unbounded complexes by [Stacks, \href{https://stacks.math.columbia.edu/tag/0F0C}{0F0C}, \href{https://stacks.math.columbia.edu/tag/0GL5}{0GL5}].

% object: proposition-adjunction-projective-line (blueprint/chap-Constructibility-of-f/sec-Duality-for-the-projective-line/proposition-adjunction-projective-line/_section.md)
\phantomsection\label{bp22}%
\begin{bpstatementit}{Proposition~6.7}{ (The adjunction for the projective line)}
Let \(S\) be of finite type over \(k\), and let \(p:\mathbf P^1_S\to S\) be the projection. The functors \(Rp_*:D((\mathbf P^1_S)_{\acute{e}t},A)\to
D(S_{\acute{e}t},A)\) and \(p^*(1)[2]:D(S_{\acute{e}t},A)\to
D((\mathbf P^1_S)_{\acute{e}t},A)\) form an adjunction

\[
Rp_* \ \dashv\ p^*(1)[2],
\tag{6.9}
\]

with unit (6.7) and counit (6.8). For every morphism \(g:S'\to S\) with \(S'\) of finite type over \(k\), the unit and counit are compatible with base change along \(g\). They are also compatible with change of finite coefficients.
\end{bpstatementit}

% proof: proof/proposition-adjunction-projective-line
\begin{proof}[{Proof}]\phantomsection\label{bp23}%
The projection formula and proper base change reduce each of the two triangle identities to the identity \(s_i=1\) of Lemma \hyperref[bp19]{6.6}, tensored with the identity of a coefficient complex.

\emph{First triangle identity.} This is the equation \(\varepsilon_{Rp_*M}\circ Rp_*(\eta_M)=1_{Rp_*M}\). Using composition of direct images, \(pp_1=pp_2\), proper base change and the projection formula, this composite equals

\[
\begin{split}
Rp_*M
&=R(pp_1)_*\Delta_*M\\
&\longrightarrow R(pp_1)_*p_1^*M(1)[2]\\
&\xrightarrow{\ \sim\ }
Rp_*\bigl(M\otimes_A^{\mathbf L}
Rp_{1*}A_{\mathbf P^1_S\times_S\mathbf P^1_S}(1)[2]\bigr)\\
&\xrightarrow{Rp_*(1_M\otimes \operatorname{Tr}_{p_1})}Rp_*M .
\end{split}
\tag{6.10}
\]

Here the first arrow is \(Rp_*\) of the first arrow of (6.7), with \(R(pp_2)_*\) rewritten as \(R(pp_1)_*\). The middle isomorphism is the projection formula for \(p_1\). It replaces the proper base change isomorphism of (6.7), after \(Rp_*\), followed by the projection formula isomorphism of (6.8). The last arrow replaces \(1\otimes\operatorname{Tr}_p\) in (6.8) by \(\operatorname{Tr}_{p_1}\), the base change of \(\operatorname{Tr}_p\) along \(p\) (Lemma \hyperref[bp14]{6.4}). In this form, the trace is applied to \(Rp_{1*}\delta\), that is, to the factor \(\Delta_*A\), and \(M\) only appears as a coefficient.

This rewriting is an equality of maps because the projection formula isomorphism is compatible with proper base change. To see this, recall that the proper base change map is the adjoint transpose of the counit of \(g^*\dashv Rg_*\) for the relevant pullback \(g\). The projection map is the adjoint transpose of the same counit tensored with the identity, and the counit of a composite adjunction is the composite of the counits. Therefore the two routes from \(R(pp_2)_*p_1^*M=R(pp_1)_*p_1^*M\) to \(Rp_*M\otimes Rp_*A(1)[2]\) have the same adjoint transpose, as in the pasting Lemma \bpout{4.5}.

On the route through (6.10) we identify \(Rp_{1*}A=p^*Rp_*A\) by proper base change and then use the projection formula for \(p\). By naturality of the projection formula in its second variable, the first arrow of (6.10) is then \(1_M\otimes Rp_{1*}\delta\) under the projection identification. The composite is therefore \(Rp_*(1_M\otimes s_1)\), which is \(1_{Rp_*M}\) by Lemma \hyperref[bp19]{6.6}.

\emph{Second triangle identity.} Here the roles of the two projections are exchanged: we must show that \(p^*(\varepsilon_N)(1)[2]\circ\eta_{p^*N(1)[2]}\) is the identity. Take \(M=p^*N(1)[2]\) in (6.7). Since \(N\) and the Tate twist are pulled back from \(S\), both \(p_1^*M\) and \(p_2^*M\) are \((pp_1)^*N(1)[2]=(pp_2)^*N(1)[2]\). The projection formula moves this common factor outside \(Rp_{2*}\), and \(\eta_{p^*N(1)[2]}\) becomes \(1\otimes Rp_{2*}\delta\). By the base change compatibility of (6.5), \(p^*(\varepsilon_N)(1)[2]\) becomes \(1\otimes\operatorname{Tr}_{p_2}\), applied to the target of \(Rp_{2*}\delta\). We obtain

\[
p^*(\varepsilon_N)(1)[2]\circ\eta_{p^*N(1)[2]}
=1_{p^*N(1)[2]}\otimes
(\operatorname{Tr}_{p_2}\circ Rp_{2*}\delta)
=1_{p^*N(1)[2]}\otimes s_2=1_{p^*N(1)[2]}.
\tag{6.11}
\]

All shifts moved in (6.10) and (6.11) are even, so no Koszul signs appear.

The unit is constructed from \(\delta\) and proper base change, and the counit from (6.5) and the projection formula. All of these ingredients are compatible with base change and change of coefficients, and hence so are \(\eta\) and \(\varepsilon\). Transported by \(\nu^*\), the same then holds for the corresponding adjunction on the pro-étale categories.
\end{proof}


% object: sec-Smooth-morphisms (blueprint/chap-Constructibility-of-f/sec-Smooth-morphisms/_section.md), content hidden
\subsection*{\texorpdfstring{6.3\quad Smooth morphisms}{6.3 Smooth morphisms}}
\phantomsection\label{bp24}%
\addcontentsline{toc}{subsection}{\texorpdfstring{6.3\quad Smooth morphisms}{6.3 Smooth morphisms}}
% text: prose-06-09
\phantomsection\label{bp25}%
Let \(f:X\to S\) be smooth of relative dimension \(d\). Extending the formula \(p^!=p^*(1)[2]\) for \(p:\mathbf P^1_S\to S\), we show in Proposition \hyperref[bp30]{6.9} that \(f^!N\simeq f^*N\otimes^{\mathbf L}\omega_f\) with \(\omega_f\) locally isomorphic to \(A(d)[2d]\). We need this formula in the proof of generic constructibility below, because it must control \(Rf_*L\) for a locally constant sheaf \(L\). Indeed, we know nothing about \(Rf_*L\) directly, whereas the finiteness theorem and proper base change control \(Rf_!\). With the formula, however, the adjunction \(Rf_!\dashv f^!\) becomes a duality between \(Rf_*L\) and \(Rf_!\) of the twisted dual of \(L\) (Lemma \hyperref[bp33]{6.10}). On a manifold, Poincaré duality exchanges ordinary and compactly supported cohomology in the same way.

For example, let \(f:\mathbf G_m\to\operatorname{Spec}\bar k\) and \(L=A=\Fell\). The localization triangle for the origin of \(\mathbf A^1\) and \(R\Gamma_c(\mathbf A^1_{\bar k},A)\simeq A(-1)[-2]\) give \(R\Gamma_c(\mathbf G_m,A)\simeq\fib\bigl(A(-1)[-2]\to A\bigr)\). Hence \(H^1_c=A\) and \(H^2_c=A(-1)\) are the only nonzero groups, and \(R\Gamma_c(\mathbf G_m,A(1)[2])\) has cohomology \(A(1)\) in degree \(-1\) and \(A\) in degree \(0\). Its dual has \(A\) in degree \(0\) and \(A(-1)\) in degree \(1\). This is indeed \(R\Gamma(\mathbf G_m,A)\), with \(H^1=A(-1)\) given by the Kummer torsors.

The duality must also commute with base change. This requires the base change isomorphism \(g'^*f^!\simeq f'^!g^*\) for smooth \(f\), and for it the local form \(\omega_f|_U\simeq A_U(d)[2d]\) on charts suffices. A global identification of \(\omega_f\) with \(A_X(d)[2d]\), by contrast, needs a trace map for \(f\); it is part of smooth purity (Chapter \bpout{9}). For the local form we use that, Zariski locally on \(X\) and \(S\), a smooth morphism \(f:X\to S\) factors as an étale morphism to some \(\mathbf A^d_S\) followed by the projection to \(S\), and the next lemma treats these two kinds of morphisms.

% object: lemma-affine-space-etale-localization (blueprint/chap-Constructibility-of-f/sec-Smooth-morphisms/lemma-affine-space-etale-localization/_section.md)
\phantomsection\label{bp26}%
\begin{bpstatementit}{Lemma~6.8}{ (Affine space and étale localization)}
Let \(S\) be of finite type over the base field \(k\).

\begin{enumerate}
\item For the projection \(p_d:\mathbf A^d_S\to S\), the right adjoint \(p_d^!\) of \(Rp_{d!}:D((\mathbf A^d_S)_{\acute{e}t},A)
\to D(S_{\acute{e}t},A)\) is \(p_d^*(d)[2d]\). Its counit \(Rp_{d!}p_d^*N(d)[2d]\to N\) is compatible with base change: for a morphism \(g:S'\to S\) of finite type over \(k\), with base change \(p'_d:\mathbf A^d_{S'}\to S'\), the pullback \(g^*\) of the counit for \(N\) is, under the base change isomorphism for \(Rp_{d!}\), the counit for \(g^*N\). The same holds for the unit, which is determined by the counit through the triangle identities.
\item Let \(f:X\to S\) be separated of finite type, let \(N\in D(S_{\acute{e}t},A)\), and let \(u:U\to X\) be separated, étale and of finite type. There is a canonical isomorphism \[
   u^*f^!N\simeq(fu)^!N
   \tag{6.12}
   \] under which the counit of \(R(fu)_!\dashv(fu)^!\) is the composite \[
   R(fu)_!u^*f^!N\simeq Rf_!u_!u^*f^!N\longrightarrow Rf_!f^!N
   \longrightarrow N
   \] of the composition isomorphism, the counit of \(u_!\dashv u^*\) and the counit of \(Rf_!\dashv f^!\). For a second separated étale morphism \(u':U'\to U\) of finite type, the isomorphism (6.12) for \(uu'\) is the composite of \(u'^*\) of (6.12) for \(u\) and (6.12) for \(u'\) and \(fu\).
\end{enumerate}
\end{bpstatementit}

% text: prose-06-10
\phantomsection\label{bp27}%
Part (2) generalizes \(j^!=j^*\) from open immersions to étale maps. For \(d=1\), \(k\) algebraically closed and \(S=\operatorname{Spec}k\), part (1) says that \(\Hom(R\Gamma_c(\mathbf A^1,F),M)=\Hom(F,M_{\mathbf A^1}(1)[2])\). For \(F=A\) both sides are \(\Hom(A,M(1)[2])\), consistent with the example above: the left side because \(R\Gamma_c(\mathbf A^1_{\bar k},A)\simeq A(-1)[-2]\), and the right side because \(R\Gamma(\mathbf A^1_{\bar k},A)=A\).

% proof: proof/lemma-affine-space-etale-localization
\begin{proof}[{Proof of Lemma 6.8}]\phantomsection\label{bp28}%
Both parts follow formally from uniqueness of right adjoints. Part (1) comes from (6.9) and \(j_!\dashv j^*\) for \(\mathbf A^1\subset\mathbf P^1\), and part (2) from \(u_!\dashv u^*\) for étale \(u\).

(1) Let \(j:\mathbf A^1_S\hookrightarrow\mathbf P^1_S\) be the open immersion, so that \(Rp_{1!}=Rp_*j_!\) for \(d=1\). Since \(j_!\dashv j^*\) and \(Rp_*\dashv p^*(1)[2]\) by (6.9), composing right adjoints gives \(p_1^!=j^*p^*(1)[2]\). This equals \(p_1^*(1)[2]\) because \(p\circ j=p_1\). For general \(d\), write \(p_d\) as a composite of \(d\) projections \(r_1,\dots,r_d\) of affine lines. By the composition isomorphism for \(Rf_!\) (Proposition proposition-composition-base-change-rf-shriek) we have \(Rp_{d!}\simeq Rr_{1!}\circ\dots\circ Rr_{d!}\). Hence \(p_d^!\) is the composite of the right adjoints of the \(Rr_{i!}\), which is \(p_d^*(d)[2d]\) by the case \(d=1\). Under this isomorphism, the counit of \(Rp_{d!}\dashv p_d^*(d)[2d]\) is the composite of the \(d\) counits for the affine lines. Each of these is a trace (6.5) composed with the counit of \(j_!\dashv j^*\). By the associativity in Proposition proposition-composition-base-change-rf-shriek, the composite does not depend on how the \(d\) factors are grouped. Each factor and the composition isomorphism commute with base change (Proposition proposition-composition-base-change-rf-shriek), and therefore so does the counit. The unit is determined by the counit through the triangle identities, and it commutes with base change as well.

(2) For an étale \(u\), the functor \(Ru_!\) is extension by zero. By its elementary construction [Stacks, \href{https://stacks.math.columbia.edu/tag/03S2}{03S2}], it is an exact functor \(u_!:\operatorname{Mod}(U_{\acute{e}t},A)\to\operatorname{Mod}(X_{\acute{e}t},A)\) left adjoint to \(u^*\). As both functors are exact, they need not be derived, and the adjunction extends to derived categories. For every \(M\in D(U_{\acute{e}t},A)\),

\[
\Hom(M,u^*f^!N)
=\Hom(u_!M,f^!N)
=\Hom(Rf_!u_!M,N)
=\Hom(R(fu)_!M,N),
\]

where the last step is the composition isomorphism of Proposition proposition-composition-base-change-rf-shriek. The Yoneda lemma therefore gives (6.12). Taking \(M=u^*f^!N\) and the identity on the left, we see that the counit of \(R(fu)_!\dashv(fu)^!\) is the composite of the counits as stated. For the cocycle condition, the isomorphism for \(uu'\) and the composite of those for \(u\) and \(u'\) both correspond under the Yoneda lemma to the bijections obtained from the composition isomorphism for \(f\circ u\circ u'\). These bijections agree by associativity of the composition isomorphisms (Proposition proposition-composition-base-change-rf-shriek).
\end{proof}

% text: prose-06-11
\phantomsection\label{bp29}%
Applied along local charts \(U\xrightarrow{e}\mathbf A^d_S\to S\), the two parts of the lemma compute \(f^!\) locally on \(X\) for every smooth \(f\). The proposition below turns this into two global statements, a tensor formula and base change.

% object: proposition-exceptional-inverse-image-smooth-morphism (blueprint/chap-Constructibility-of-f/sec-Smooth-morphisms/proposition-exceptional-inverse-image-smooth-morphism/_section.md)
\phantomsection\label{bp30}%
\begin{bpstatementit}{Proposition~6.9}{ (The exceptional inverse image of a smooth morphism)}
Let \(f:X\to S\) be separated, smooth and of finite type, of pure relative dimension \(d\), and let \(\omega_f=f^!A_S\in D(X_{\acute{e}t},A)\) be its dualizing complex. Let \(N\in D(S_{\acute{e}t},A)\).

\begin{enumerate}
\item On an open \(U\subset X\) with \(f|_U=p_d\circ e\), where \(e:U\to\mathbf A^d_S\) is separated and étale, we have \((f^!N)|_U\simeq e^*p_d^*N(d)[2d]\), with counit the sum over the fibres of \(e\) followed by the counit of Lemma lemma-affine-space-etale-localization. In particular \(\omega_f|_U\simeq A_U(d)[2d]\), and \(\omega_f[-2d]\) is a locally constant sheaf of free \(A\)-modules of rank one.
\item The canonical map \[
   \mu_N:f^*N\otimes_A^{\mathbf L}\omega_f\longrightarrow f^!N
   \tag{6.13}
   \] is an isomorphism.
\item Let \(g:S'\to S\) be a morphism, let \(X'=X\times_SS'\), and let \(f':X'\to S'\), \(g':X'\to X\) be the projections. Then the base change map \(\beta_N:g'^*f^!N\to f'^!g^*N\) of Definition \hyperref[bp7]{6.2} is an isomorphism. It identifies \(g'^*\omega_f\) with \(\omega_{f'}\) and is compatible with \(\mu_N\).
\end{enumerate}
\end{bpstatementit}

% proof: proof/proposition-exceptional-inverse-image-smooth-morphism
\begin{proof}[{Proof}]\phantomsection\label{bp31}%
The maps \(\mu_N\) and \(\beta_N\) are defined globally as adjoint transposes. Since being an isomorphism is local on \(X\), respectively on \(X'\), we show that on a chart both maps become identifications of inverse images under the explicit identification (6.14). For this we follow the counits through the factorization \(U\to\mathbf A^d_V\to V\).

We write \(\varepsilon_N:Rf_!f^!N\to N\) for the counit of \(Rf_!\dashv f^!\) (for \(f=p\) this is (6.8)). The map \(\mu_N\) is the transpose of the morphism

\[
Rf_!(f^*N\otimes^{\mathbf L}\omega_f)
\xrightarrow{\ \sim\ }
N\otimes^{\mathbf L}Rf_!\omega_f
\xrightarrow{1\otimes\varepsilon_{A_S}}N ,
\]

where the first arrow is the projection formula. The map \(\beta_N\) is the transpose, for \(f'\), of the morphism (6.2) (Definition \hyperref[bp7]{6.2}).

\emph{Local charts.} The computation of \(f^!\) on a chart also proves (1). Zariski locally on \(X\) and \(S\), a smooth morphism is an étale map to affine space followed by the projection [Stacks, \href{https://stacks.math.columbia.edu/tag/054L}{054L}]. Thus every point of \(X\) has an affine open neighbourhood \(u:U\hookrightarrow X\) mapping into an affine open \(v:V\hookrightarrow S\), together with a separated étale map \(e:U\to\mathbf A^d_V\) such that \(fu=vp_de\), where \(p_d:\mathbf A^d_V\to V\) is the projection. By Lemma \hyperref[bp26]{6.8} and uniqueness of right adjoints, we have

\[
u^*f^!N=e^*p_d^*v^*N(d)[2d],\qquad
u^*\omega_f=A_U(d)[2d].
\tag{6.14}
\]

In particular, \(\omega_f\) has a single nonzero cohomology sheaf, in degree \(-2d\), which is locally free of rank one. This proves (1), because the counit on \(U\) is the composite of the étale counit for \(e\) and the counit for \(p_d\). The first of these is the sum over the fibres.

\emph{\(\mu_N\) is an isomorphism.} On each such \(U\), under (6.14), the source and the target of \(\mu_N\) are both \(e^*p_d^*v^*N(d)[2d]\). We claim that the map is the identification obtained by reassociating tensor product and twist. It suffices to compare the transposes. Both are the projection formula followed by the counit for \(v\), the \(d\) counits for the affine lines, and the étale counit for \(e\). They agree by the formula for the counit of a composite in Lemma \hyperref[bp26]{6.8}. Hence \(\mu_N\) is an isomorphism on \(X\).

\emph{\(\beta_N\) is an isomorphism.} For each type of chart, under the explicit description of \(f^!\) as a twisted inverse image, \(\beta_N\) is the evident identification of inverse images. For the projective line, the transpose of (6.2) is the usual identification \(g'^*p^*N(1)[2]=p'^*g^*N(1)[2]\) of pullbacks. This holds because the unit (6.7) and counit (6.8) commute with base change. By restriction to the affine line and composition, the same holds for affine space. For a separated étale map \(u\), the transpose is again the usual identification of inverse images. Indeed, the stalk of \(u_!M\) at a geometric point is the direct sum of the stalks of \(M\) at the geometric points over it. This description is compatible with base change, and under it the counit \(u_!u^*N\to N\) is the sum of the identity maps of the stalks of \(N\).

The cases combine, because base change for \(Rf_!\) is compatible with composition (Proposition proposition-composition-base-change-rf-shriek) and the counit of a composite is the composite of the counits. The open immersions \(u\) and \(v\) are étale, so the étale case covers them as well. Substituting these facts into (6.2), we see that the restriction of \(\beta_N\) to \(U\times_XX'\) is the composite of the étale and affine-space cases just treated. Hence \(\beta_N\) is an isomorphism on a cover of \(X'\), and therefore it is an isomorphism.

\emph{Compatibilities.} Taking \(N=A_S\), we see that \(\beta_N\) identifies \(g'^*\omega_f\) with \(\omega_{f'}\), because \(g^*A_S=A_{S'}\). The compatibility of \(\mu_N\) with \(\beta_N\) says that two maps \(g'^*(f^*N\otimes^{\mathbf L}\omega_f)\to f'^!g^*N\) agree. Here the first is \(\mu_N\) followed by \(\beta_N\), and the second is the identification \(g'^*(f^*N\otimes^{\mathbf L}\omega_f)=f'^*g^*N\otimes^{\mathbf L}\omega_{f'}\), coming from \(g'^*f^*=f'^*g^*\) and the case \(N=A_S\), followed by \(\mu_N\) for \(f'\). Their transposes are maps \(Rf'_!g'^*(f^*N\otimes^{\mathbf L}\omega_f)\to g^*N\). Both are given by the projection formula followed by the counit, once before and once after the base change isomorphism for \(Rf_!\). They agree for two reasons. First, the projection formula isomorphism is compatible with the base change isomorphism for \(Rf_!\). Second, by the definition (6.2) of \(\beta_N\), the counit for \(f'\) composed with \(Rf'_!\) of \(\beta_N\) is \(g^*\) of the counit for \(f\) under that isomorphism.
\end{proof}

% text: prose-06-12
\phantomsection\label{bp32}%
In Step 2 of the proof of Theorem \hyperref[bp51]{6.17}, the next lemma is applied to \(f=bj\), with \(b:Y\to S\) projective and \(j:X\hookrightarrow Y\) open.

% object: lemma-smooth-morphisms-locally-constant-sheaves (blueprint/chap-Constructibility-of-f/sec-Smooth-morphisms/lemma-smooth-morphisms-locally-constant-sheaves/_section.md)
\phantomsection\label{bp33}%
\begin{bpstatementit}{Lemma~6.10}{ (Smooth morphisms and locally constant sheaves)}
Let \(S\) be integral and let \(f:X\to S\) be separated, smooth and of finite type. For a finite locally constant \(\Fell\)-sheaf \(L\) on \(X_{\acute{e}t}\), there is a nonempty open \(S_0\subset S\) with the following property. For the restriction \(f_{S_0}:X_{S_0}\to S_0\), the complex \(Rf_{S_0,*}L\) is bounded with finite locally constant cohomology sheaves, and its formation commutes with base change. Namely, for every \(k\)-scheme \(v:T\to S_0\) of finite type, with base change \(f_T:X_T\to T\) and \(v_X:X_T\to X_{S_0}\), the canonical base change map \(v^*Rf_{S_0,*}L\to Rf_{T,*}v_X^*L\) is an isomorphism.
\end{bpstatementit}

% proof: proof/lemma-smooth-morphisms-locally-constant-sheaves
\begin{proof}[{Proof}]\phantomsection\label{bp34}%
We write \(Rf_*L\) as the dual of \(Rf_!N\) for a suitable constructible \(N\). The point is that constructibility and base change are known for \(Rf_!N\), and over a dense open of \(S\) the cohomology sheaves of \(Rf_!N\) are locally constant, and there dualizing preserves both properties.

By Proposition \bpout{2.42}, the locally constant sheaf \(L\) comes from a locally constant sheaf on the étale site, and we work with the latter.

Since \(X\) is the disjoint union of finitely many open pieces of pure relative dimension, we may assume that \(f\) has pure relative dimension \(d\). Let \(\omega_f=f^!\mathbf F_{\ell,S}\) as in Proposition \hyperref[bp30]{6.9}, and let

\[
N=\cHom_X(L,\mathbf F_{\ell,X})\otimes^{\mathbf L}\omega_f .
\]

By Proposition \hyperref[bp30]{6.9}, \(\omega_f\) is locally \(\mathbf F_{\ell,X}(d)[2d]\). Hence \(N\) is bounded constructible and \(\cHom_X(N,\omega_f)=L\).

\emph{Duality.} Consider the pairing

\[
Rf_!N\otimes^{\mathbf L}Rf_*L
\longrightarrow Rf_!(N\otimes^{\mathbf L}L)
\longrightarrow Rf_!\omega_f
\xrightarrow{\varepsilon_{\mathbf F_{\ell,S}}}\mathbf F_{\ell,S},
\tag{6.15}
\]

where the first arrow is the projection formula followed by the counit \(f^*Rf_*L\to L\), and the second is evaluation. Its adjoint is a map

\[
\delta_L:Rf_*L
\longrightarrow \cHom_S(Rf_!N,\mathbf F_{\ell,S}),
\tag{6.16}
\]

and we claim that it is an isomorphism. Since the map is an adjoint transpose, we test it on Hom groups out of a family of generators, where adjunctions turn one side into the other. To this end we write \(L=\cHom_X(N,f^!\mathbf F_{\ell,S})\), and test \(\delta_L\) on the objects \(P[n]\), where \(P=u_!\Fell\) and \(u:V\to S\) is quasi-compact, separated and étale. The objects \(u_!\Fell\) are constructible of finite Tor amplitude, and their shifts detect isomorphisms in \(D(S_{\acute{e}t},\Fell)\). Since \(f^*P\otimes^{\mathbf L}N\) lies in \(D^b_{\mathrm{ctf}}(X,\Fell)\), we may use tensor-Hom adjunction, the adjunction \(Rf_!\dashv f^!\) of Theorem \bpout{5.31} between \(\Loc(D^b_{\mathrm{ctf}}(X,\Fell))\) and \(\Loc(D^b_{\mathrm{ctf}}(S,\Fell))\), and the projection formula. With Hom groups in \(D(S_{\acute{e}t},\Fell)\), we then obtain

\[
\Hom(P,Rf_*\cHom(N,f^!\mathbf F_{\ell,S}))
=\Hom(Rf_!(f^*P\otimes^{\mathbf L}N),\mathbf F_{\ell,S})
=\Hom(P\otimes^{\mathbf L}Rf_!N,\mathbf F_{\ell,S}).
\]

By tensor-Hom adjunction, the last group is \(\Hom(P,\cHom_S(Rf_!N,\mathbf F_{\ell,S}))\). Unwinding the adjunctions, one sees that the composite bijection is the map induced by \(\delta_L\), for every shift of \(P\). Hence \(\delta_L\) is an isomorphism.

\emph{Shrinking the base.} Since \(\delta_L\) is an isomorphism, it now suffices to understand the dual of \(Rf_!N\). By the finiteness theorem for \(Rf_!\) [Stacks, \href{https://stacks.math.columbia.edu/tag/0GL0}{0GL0}, \href{https://stacks.math.columbia.edu/tag/0G2A}{0G2A}], the complex \(Rf_!N\) is bounded constructible. A constructible sheaf on the integral scheme \(S\) is locally constant on a nonempty open subset. Thus each of the finitely many nonzero cohomology sheaves of \(Rf_!N\) is locally constant on a nonempty open of \(S\). Let \(S_0\) be the intersection of these opens; it is nonempty because \(S\) is irreducible.

Over the field \(\Fell\), a finite locally constant sheaf is locally free of finite rank. Its dual is exact, and evaluation identifies it with its double dual. We now argue by induction on the number of nonzero cohomology sheaves, using that the relevant canonical maps form morphisms of the truncation triangles. In this way we see that a bounded complex with finite locally constant cohomology is dualizable, that its dual is again of this kind, and that forming the dual commutes with every pullback. Through \(\delta_L\), this proves the first assertion over \(S_0\).

\emph{Base change.} The duality \(\delta_L\) exchanges \(Rf_*L\) with the dual of \(Rf_!N\), and \(Rf_!N\) satisfies base change. We show that the duality is compatible with base change; then \(Rf_*L\) satisfies base change too. Let \(v:T\to S_0\) be a \(k\)-scheme of finite type over \(S_0\), and let \(f_T:X_T\to T\) and \(v_X:X_T\to X\) be the base changes. First, base change for \(Rf_!\) and (6.1) identify \(v_X^*N\), \(v^*Rf_!N\) and \(v_X^*\omega_f\) with the corresponding objects for \(f_T\). Next, the base change map \(v^*Rf_*L\to Rf_{T,*}v_X^*L\) is characterized by the fact that it carries the pullback of the counit \(f^*Rf_*L\to L\) to the counit for \(f_T\). Hence the pullback of (6.15) is the pairing for \(f_T\), composed with this base change map. Finally, the complex \(v^*Rf_!N\) is dualizable with dual \(v^*\cHom(Rf_!N,\mathbf F_{\ell,S})\). Therefore the base change map for \(Rf_*L\) is identified with the dual of the base change isomorphism for \(Rf_!N\), and it is an isomorphism.
\end{proof}

% object: example-shrinking-base-smooth (blueprint/chap-Constructibility-of-f/sec-Smooth-morphisms/example-shrinking-base-smooth/_section.md)
\phantomsection\label{bp35}%
\begin{bpstatement}{Example~6.11}{ (Why the base must be shrunk)}
Let \(k\) be algebraically closed, \(S=\mathbf A^1\) with coordinate \(t\), and let \(X\subset\mathbf A^1_S\) be the complement of the two sections \(x=0\) and \(x=t\), with \(f:X\to S\) the projection and \(L=\Fell\). The morphism \(f\) is smooth. For \(t\neq0\) its fibre is \(\mathbf A^1\) minus two points, that is, a projective line minus three points. Its \(H^1\) is \(\Fell(-1)^{\oplus2}\). The fibre over \(t=0\), however, is \(\mathbf G_m\), with \(H^1=\Fell(-1)\). On \(S_0=\mathbf G_m\) the lemma applies, but no open containing the origin works. Suppose \(R^1f_*L\) were locally constant near \(0\) and commuted with base change to the points of \(S\). Then its stalks at \(0\) and at nearby \(t\) would be the \(H^1\) of the fibres, which have different ranks. What happens is that the two missing points collide at \(t=0\), and \(Rf_*\) cannot see this from the generic fibre.
\end{bpstatement}


% object: sec-Coordinate-projections (blueprint/chap-Constructibility-of-f/sec-Coordinate-projections/_section.md), content hidden
\subsection*{\texorpdfstring{6.4\quad Coordinate projections}{6.4 Coordinate projections}}
\phantomsection\label{bp36}%
\addcontentsline{toc}{subsection}{\texorpdfstring{6.4\quad Coordinate projections}{6.4 Coordinate projections}}
% text: prose-06-13
\phantomsection\label{bp37}%
To control \(Rj_*F\) for an open immersion \(j:X\hookrightarrow Y\) over \(S\), with \(Y\) projective over \(S\), we cover most of \(Y\) by opens where the induction hypothesis applies and treat the rest separately. On a chart \(Y_\alpha\), a coordinate function \(t:Y_\alpha\to\mathbf A^1_S\) lets us view \(X\cap Y_\alpha\) over the new base \(\mathbf A^1_S\), over which its generic fibre has one dimension less. Hence the induction hypothesis in Theorem \hyperref[bp51]{6.17} gives constructibility and base change over a dense open of \(\mathbf A^1_S\). The preimages of these opens, however, do not cover \(Y\). Take for instance \(Y=\mathbf P^1_S\) with the two standard charts: there a point is missed only if each of the coordinates \(x\) and \(1/x\) defined there takes an excluded value. After shrinking \(S\), the leftover set is therefore a closed subset finite over \(S\). Lemma \hyperref[bp41]{6.13} shows that the same holds in general.

On the leftover finite part we cannot argue by dimension again. We do not have to, since a finite direct image is exact and conservative. This means that on a subset finite over \(S\), the constructibility of a sheaf can be read off from its direct image to \(S\). Lemma \hyperref[bp38]{6.12} turns this into a criterion. If a complex is already constructible away from such a subset, a proper direct image to \(S\) detects whether it is constructible everywhere. In the smooth case this direct image is \(R(bj)_*F\), which Lemma \hyperref[bp33]{6.10} controls. For this reason the induction passes through the smooth case first.

% object: lemma-detection-proper-direct-image (blueprint/chap-Constructibility-of-f/sec-Coordinate-projections/lemma-detection-proper-direct-image/_section.md)
\phantomsection\label{bp38}%
\begin{bpstatementit}{Lemma~6.12}{ (Detection by a proper direct image)}
Let \(b:Y\to S\) be proper, let \(j:U\hookrightarrow Y\) be an open immersion with reduced closed complement \(i:Z\hookrightarrow Y\), and suppose that \(bi:Z\to S\) is finite.

\begin{enumerate}
\item If \(K\in D^+(Y_{\acute{e}t},\Fell)\), \(j^*K\in\Dbc(U,\Fell)\), and \(Rb_*K\in\Dbc(S,\Fell)\), then \(K\in\Dbc(Y,\Fell)\).
\item If \(\varphi:K\to K'\) is a morphism in \(D^+(Y_{\acute{e}t},\Fell)\) such that \(j^*\varphi\) and \(Rb_*\varphi\) are isomorphisms, then \(\varphi\) is an isomorphism.
\end{enumerate}
\end{bpstatementit}

% proof: proof/lemma-detection-proper-direct-image
\begin{proof}[{Proof}]\phantomsection\label{bp39}%
Using the localization triangle we split \(K\) into a part on \(U\), which is known, and a part on \(Z\). Since \(Z\) is finite over \(S\), the part on \(Z\) is detected by its direct image to \(S\). We compute this direct image from \(Rb_*K\) and the known part.

We work on the étale sites, in \(D^+(-_{\acute{e}t},\Fell)\) and \(\operatorname{Mod}(-_{\acute{e}t},\Fell)\).

\emph{Finite direct images.} Let \(h:Z\to S\) be finite. Then \(h_*\) is exact and conservative, and \(Rh_*M\in\Dbc\) implies \(M\in\Dbc\). The first two properties hold because, for a sheaf \(F\) of \(\Fell\)-modules on \(Z\), \(h_*F\) has stalks

\[
(h_*F)_{\bar s}
=\prod_{\bar z\in |Z_{\bar s}|}F_{\bar z}
\tag{6.17}
\]

by [Stacks, \href{https://stacks.math.columbia.edu/tag/03QP}{03QP}]. Suppose \(h_*F\) is constructible. The counit \(h^*h_*F\to F\) is surjective, since on stalks it is a projection from the finite product (6.17). Its source is constructible, and a quotient of a constructible sheaf is constructible [Stacks, \href{https://stacks.math.columbia.edu/tag/09BH}{09BH}]. Hence \(F\) is constructible. For a complex \(M\) we have \(H^n(Rh_*M)=h_*H^n(M)\) because \(h_*\) is exact, so the statement for complexes follows from conservativity and the sheaf case in each degree.

\emph{The localization triangle.} The exact sequence for an open and its closed complement [Stacks, \href{https://stacks.math.columbia.edu/tag/095L}{095L}], applied termwise to \(K\), gives a triangle

\[
j_!j^*K\longrightarrow K\longrightarrow i_*i^*K
\longrightarrow j_!j^*K[1].
\tag{6.18}
\]

\emph{Part (1).} The first term \(j_!j^*K\) is bounded constructible, because \(j^*K\) is and extension by zero preserves constructibility. Its direct image \(Rb_*j_!j^*K\) is bounded constructible by the finiteness theorem [Stacks, \href{https://stacks.math.columbia.edu/tag/0GL0}{0GL0}, \href{https://stacks.math.columbia.edu/tag/0G2A}{0G2A}], which applies since \(Rb_*=Rb_!\) for proper \(b\). Applying \(Rb_*\) to (6.18), we obtain a triangle whose first two terms, \(Rb_*j_!j^*K\) and \(Rb_*K\), are bounded constructible. Hence the third term \(Rb_*i_*i^*K=R(bi)_*i^*K\) is bounded constructible too. The first step, applied to \(h=bi\), now gives \(i^*K\in\Dbc(Z,\Fell)\). As \(i_*\) is exact and preserves constructibility, (6.18) shows that \(K\in\Dbc(Y,\Fell)\).

\emph{Part (2).} Let \(M\) be a cone of \(\varphi\); we must show \(M=0\). We have \(j^*M=0\) because \(j^*\varphi\) is an isomorphism. By (6.18), this means \(M=i_*i^*M\). Because \(Rb_*\varphi\) is an isomorphism, we also have \(0=Rb_*M=R(bi)_*i^*M\). Since \((bi)_*\) is exact, the cohomology sheaves of \(R(bi)_*i^*M\) are the sheaves \((bi)_*H^n(i^*M)\). Conservativity of \((bi)_*\) then shows that every \(H^n(i^*M)\) vanishes, and so \(M=i_*i^*M=0\).
\end{proof}

% text: prose-06-14
\phantomsection\label{bp40}%
The next lemma produces such a finite complement. Projections to coordinate lines lower the dimension of the generic fibre, and the subset that they do not cover becomes finite over a smaller base.

% object: lemma-coordinate-projections (blueprint/chap-Constructibility-of-f/sec-Coordinate-projections/lemma-coordinate-projections/_section.md)
\phantomsection\label{bp41}%
\begin{bpstatementit}{Lemma~6.13}{ (Coordinate projections)}
Let \(S\) be an integral Noetherian affine scheme with function field \(k(S)\), and let \(Y\to S\) be projective. Choose a finite affine open cover \(Y=\bigcup_\alpha Y_\alpha\) and closed immersions \(Y_\alpha\hookrightarrow\mathbf A^{m_\alpha}_S\) with \(m_\alpha\geq1\), and write \(t_{\alpha r}:Y_\alpha\to\mathbf A^1_S\) (\(1\le r\le m_\alpha\)) for the coordinate functions.

\begin{enumerate}
\item If \(X\subset Y\) is open and \(\dim X_{k(S)}\leq n\), then for every \(\alpha,r\) the generic fibre of \(t_{\alpha r}:X\cap Y_\alpha\to\mathbf A^1_S\) has dimension at most \(n-1\).
\item Suppose given a dense open \(V_{\alpha r}\subset\mathbf A^1_S\) for each \(\alpha,r\), and let \[
   G=\bigcup_{\alpha,r}
   t_{\alpha r}^{-1}(V_{\alpha r})\subset Y.
   \] After shrinking \(S\), the reduced closed complement \(Z=Y\setminus G\) is finite over \(S\).
\end{enumerate}
\end{bpstatementit}

% proof: proof/lemma-coordinate-projections
\begin{proof}[{Proof}]\phantomsection\label{bp42}%
Part (1) is a transcendence degree count. For part (2), note that a point of \(Z\) has all its coordinates in the finite sets of excluded values. After shrinking \(S\), these sets are cut out by monic polynomials, and monic equations in every coordinate force finiteness.

(1) We use that the dimension of an integral scheme of finite type over a field is the transcendence degree of its function field. Consider an irreducible component of \(X_{k(S)}\cap Y_\alpha\) with function field \(L\). If \(t_{\alpha r}\) is algebraic over \(k(S)\) on it, the component does not meet the generic fibre of \(t_{\alpha r}\), as it does not dominate \(\mathbf A^1_{k(S)}\). Otherwise \(t_{\alpha r}\) is transcendental, and

\[
\operatorname{trdeg}_{k(S)(t_{\alpha r})}L
=\operatorname{trdeg}_{k(S)} L-1\leq n-1.
\]

The left-hand side is the dimension of the corresponding component of the generic fibre of \(t_{\alpha r}\), and every component of that generic fibre arises in this way.

(2) Let \(t\) be the coordinate of \(\mathbf A^1_S\). The ideal of the complement of \(V_{\alpha r}\) contains a nonzero polynomial \(P_{\alpha r}\in\mathcal O(S)[t]\). Indeed, since \(V_{\alpha r}\) is dense, its complement is a proper closed subset of the integral scheme \(\mathbf A^1_S\), and its ideal is nonzero. After inverting the leading coefficients of these finitely many polynomials, we may assume that all of them are monic. The function \(P_{\alpha r}(t_{\alpha r})\) is zero in the coordinate ring of the reduced affine scheme \(Z\cap Y_\alpha\), because it vanishes at every point, \(t_{\alpha r}\) mapping \(Z\cap Y_\alpha\) into the complement of \(V_{\alpha r}\). This ring is therefore a quotient of

\[
\mathcal O(S)[t_{\alpha1},\ldots,t_{\alpha m_\alpha}]/
(P_{\alpha1}(t_{\alpha1}),\ldots,
P_{\alpha m_\alpha}(t_{\alpha m_\alpha})),
\]

which is generated as an \(\mathcal O(S)\)-module by the monomials whose exponent in each variable is less than the degree of the corresponding polynomial. Each \(Z\cap Y_\alpha\) is thus finite over \(S\), and \(Z\to S\) has finite fibres. As \(Z\) is closed in the projective scheme \(Y\), the morphism \(Z\to S\) is also proper, and a proper morphism with finite fibres is finite [Stacks, \href{https://stacks.math.columbia.edu/tag/02LS}{02LS}].
\end{proof}


% object: sec-Generic-base-change (blueprint/chap-Constructibility-of-f/sec-Generic-base-change/_section.md), content hidden
\subsection*{\texorpdfstring{6.5\quad Generic base change}{6.5 Generic base change}}
\phantomsection\label{bp43}%
\addcontentsline{toc}{subsection}{\texorpdfstring{6.5\quad Generic base change}{6.5 Generic base change}}
% text: prose-06-15
\phantomsection\label{bp44}%
Let \(j:X\hookrightarrow Y\) be an open immersion of schemes over an integral base \(S\), and let \(F\) be a constructible \(\Fell\)-sheaf on \(X\). For such data, Deligne's generic base change theorem says that over a dense open of \(S\) the complex \(Rj_*F\) is bounded constructible and its formation commutes with base change. In particular, over a field it gives the finiteness of \(Rj_*\) (Corollary \hyperref[bp54]{6.18}), and this is all that the computation of \(f^!\) requires. On the other hand, the punctured plane over the \(t\)-line, discussed at the beginning of this chapter, shows that we cannot ask for more than a dense open.

In characteristic \(p\) the generic fibre need not contain a smooth dense open. For example, let \(X=\mathbf A^1_{\mathbf F_p}\) map to the \(t\)-line by \(x\mapsto x^p\). The generic fibre is \(\operatorname{Spec}\mathbf F_p(t^{1/p})\), which is not smooth over \(\mathbf F_p(t)\). For this reason Step 3 of the proof of Theorem \hyperref[bp51]{6.17} below allows a finite universal homeomorphism of the base. Such a map changes neither the étale topoi nor the conclusion.

We fix notation for base change. With \(j\) as above and a morphism \(T\to S\), we write \(j_T:X_T\hookrightarrow Y_T\) for the base change and \(g_X:X_T\to X\), \(g_Y:Y_T\to Y\) for the projections. The base change map \(g_Y^*Rj_*K\to Rj_{T,*}g_X^*K\) is then the map (4.5).

The proof of Theorem \hyperref[bp51]{6.17} is an induction on the dimension of the generic fibre, in three steps. The next three lemmas are the parts of this proof that do not use the induction hypothesis, and each serves one step. First, Lemma \hyperref[bp45]{6.14} is Step 2: if \(X\) is smooth and \(F\) locally constant, it extends the conclusion from an open \(G\) with complement finite over \(S\), provided by Step 1, across that complement. Next, Lemma \hyperref[bp47]{6.15} allows Step 3 to replace \(S\) by a finite universal homeomorphism, as the example above requires. Lemma \hyperref[bp49]{6.16} then provides, after this replacement, the dense open \(V\) of Step 3, smooth with \(F\) locally constant, whose complement has a smaller generic fibre.

% object: lemma-extension-across-finite-subset (blueprint/chap-Constructibility-of-f/sec-Generic-base-change/lemma-extension-across-finite-subset/_section.md)
\phantomsection\label{bp45}%
\begin{bpstatementit}{Lemma~6.14}{ (Extension across a finite subset)}
Let \(S\) be integral and affine, let \(b:Y\to S\) be projective, and let \(j:X\hookrightarrow Y\) be an open immersion with \(X\) smooth over \(S\). Let \(F\) be a finite locally constant \(\Fell\)-sheaf on \(X_{\acute{e}t}\). Suppose that there is an open \(G\subset Y\) whose reduced complement \(i:Z\hookrightarrow Y\) is finite over \(S\), such that \((Rj_*F)|_G\in\Dbc(G,\Fell)\) and such that for every \(k\)-scheme \(T\to S\) of finite type the base change map \(g_Y^*Rj_*F\to Rj_{T,*}g_X^*F\) is an isomorphism on \(G_T\). Then there is a nonempty open \(S_0\subset S\) such that \(Rj_{S_0,*}(F|_{X_{S_0}})\in\Dbc(Y_{S_0},\Fell)\) and such that for every \(k\)-scheme \(T\to S_0\) of finite type the base change map \(g_Y^*Rj_{S_0,*}(F|_{X_{S_0}})\to Rj_{T,*}g_X^*F\) is an isomorphism.
\end{bpstatementit}

% proof: proof/lemma-extension-across-finite-subset
\begin{proof}[{Proof}]\phantomsection\label{bp46}%
Lemma \hyperref[bp33]{6.10} controls the direct image \(R(bj)_*F\) of \(Rj_*F\) to \(S\). Then Lemma \hyperref[bp38]{6.12} turns this, together with the hypothesis on \(G\), into control of \(Rj_*F\) along \(Z\). Throughout, we work on the étale sites.

For constructibility we check the hypotheses of part (1) of Lemma \hyperref[bp38]{6.12} for \(K=Rj_*F\) and \(U=G\). The complex \(Rj_*F\) lies in \(D^+(Y_{\acute{e}t},\Fell)\) because \(F\) is a sheaf. Its restriction to \(G\) is in \(\Dbc\) by assumption, and \(Z=Y\setminus G\) is finite over \(S\). It remains to treat the direct image. Since \(bj:X\to S\) is smooth and \(F\) is locally constant, we may apply Lemma \hyperref[bp33]{6.10}. After shrinking \(S\) to a nonempty open \(S_0\), the complex \(Rb_*Rj_*F=R(bj)_*F\) lies in \(\Dbc(S_0,\Fell)\), and its formation commutes with base change along every \(T\to S_0\) of finite type over \(k\). Part (1) of Lemma \hyperref[bp38]{6.12} then gives \(Rj_*F\in\Dbc(Y_{S_0},\Fell)\).

For base change we use part (2) of the lemma instead. Let \(T\to S_0\) be a \(k\)-scheme of finite type, and let \(\varphi\) be the base change map \(g_Y^*Rj_*F\to Rj_{T,*}g_X^*F\). By assumption it is an isomorphism on \(G_T\). Proper base change for \(b\) and the compatibility of base change maps with composition [Stacks, \href{https://stacks.math.columbia.edu/tag/0E46}{0E46}] identify \(Rb_{T,*}\varphi\) with the base change map for \(R(bj)_*F\), which is an isomorphism. Since \(b_T\) is proper and \(b_Ti_T\) is finite, part (2) of Lemma \hyperref[bp38]{6.12} applies over \(T\). It shows that \(\varphi\) is an isomorphism.
\end{proof}

% object: lemma-universal-homeomorphism-base (blueprint/chap-Constructibility-of-f/sec-Generic-base-change/lemma-universal-homeomorphism-base/_section.md)
\phantomsection\label{bp47}%
\begin{bpstatementit}{Lemma~6.15}{ (Universal homeomorphisms of the base)}
Let \(S\) be integral, let \(j:X\hookrightarrow Y\) be an open immersion of \(S\)-schemes, and let \(F\) be a constructible \(\Fell\)-sheaf on \(X_{\acute{e}t}\). Let \(S'\to S\) be a finite universal homeomorphism with \(S'\) integral, let \(X'=(X_{S'})_{\mathrm{red}}\) and \(Y'=(Y_{S'})_{\mathrm{red}}\), let \(j':X'\hookrightarrow Y'\) be the open immersion, and let \(F'\) be the pullback of \(F\) to \(X'\). Suppose that there is a nonempty open \(U'\subset S'\) such that \(Rj'_{U',*}(F'|_{X'_{U'}})\in\Dbc(Y'_{U'},\Fell)\) and such that for every \(k\)-scheme \(T'\to U'\) of finite type the base change map for \(j'\) along \(T'\) is an isomorphism. Then there is a nonempty open \(U\subset S\) with the same two properties for \(j\) and \(F\).
\end{bpstatementit}

% proof: proof/lemma-universal-homeomorphism-base
\begin{proof}[{Proof}]\phantomsection\label{bp48}%
A finite universal homeomorphism is a homeomorphism of underlying spaces, so we may remove from \(S\) the closed image of \(S'\setminus U'\). The remaining open \(U\) is nonempty, and its inverse image \(S'_U\) lies in \(U'\). Over \(U\) the morphisms \(X'_U\to X_U\) and \(Y'_U\to Y_U\) are finite universal homeomorphisms, and pullback along them is therefore an equivalence of étale topoi [Stacks, \href{https://stacks.math.columbia.edu/tag/03SI}{03SI}]. Such equivalences preserve constructibility in both directions and commute with \(Rj_*\). Hence pullback along these morphisms detects boundedness and constructibility, and we get \(Rj_{U,*}(F|_{X_U})\in\Dbc(Y_U,\Fell)\).

For base change, let \(T\to U\) be a \(k\)-scheme of finite type, with projections \(g_X:X_T\to X\) and \(g_Y:Y_T\to Y\), and let

\[
\beta_T:g_Y^*Rj_*F\longrightarrow Rj_{T,*}g_X^*F
\]

be the base change map for \(j\) along \(T\). Put \(T'=T\times_US'_U\). Then \(v:T'\to T\) is again a universal homeomorphism, and so are the induced maps \(v_X:X'_{T'}\to X_T\) and \(v_Y:Y'_{T'}\to Y_T\), each a thickening followed by a base change of \(v\). Let \(\beta'_{T'}:g_{Y'}^*Rj'_*F'\to Rj'_{T',*}g_{X'}^*F'\) be the base change map for \(j'\) along \(T'\). Base change maps compose [Stacks, \href{https://stacks.math.columbia.edu/tag/0EZR}{0EZR}, \href{https://stacks.math.columbia.edu/tag/0EZS}{0EZS}], and \(v_Y^*Rj_{T,*}\simeq Rj'_{T',*}v_X^*\) because \(v_X^*\) and \(v_Y^*\) are equivalences of étale topoi [Stacks, \href{https://stacks.math.columbia.edu/tag/03SI}{03SI}]. Hence \(v_Y^*\beta_T\) is identified with \(\beta'_{T'}\). The map \(\beta'_{T'}\) is an isomorphism by assumption, so \(v_Y^*\beta_T\) is one. Since \(v_Y^*\) is an equivalence, it is conservative, and therefore \(\beta_T\) is an isomorphism.
\end{proof}

% object: lemma-generic-smooth-open (blueprint/chap-Constructibility-of-f/sec-Generic-base-change/lemma-generic-smooth-open/_section.md)
\phantomsection\label{bp49}%
\begin{bpstatementit}{Lemma~6.16}{ (A smooth dense open after a universal homeomorphism)}
Let \(S\) be integral, let \(X\) be an \(S\)-scheme all of whose irreducible components dominate \(S\), let \(n=\dim X_{k(S)}\), and let \(F\) be a constructible \(\Fell\)-sheaf on \(X_{\acute{e}t}\). After shrinking \(S\) to a nonempty open, there are a finite universal homeomorphism \(S'\to S\) with \(S'\) integral and a dense open \(V\subset(X_{S'})_{\mathrm{red}}\) such that \(V\) is smooth over \(S'\), the pullback of \(F\) to \(V\) is locally constant, and the generic fibre of \((X_{S'})_{\mathrm{red}}\setminus V\) over \(S'\) has dimension less than \(n\).
\end{bpstatementit}

% proof: proof/lemma-generic-smooth-open
\begin{proof}[{Proof}]\phantomsection\label{bp50}%
By [Stacks, \href{https://stacks.math.columbia.edu/tag/0550}{0550}], after shrinking \(S\), there is a finite universal homeomorphism \(S'\to S\) with \(S'\) integral, such that \((X_{S'})_{\mathrm{red}}\to S'\) is flat of finite presentation with geometrically reduced generic fibre. This generic fibre has a dense smooth open [Stacks, \href{https://stacks.math.columbia.edu/tag/056V}{056V}]. By the fibrewise criterion [Stacks, \href{https://stacks.math.columbia.edu/tag/01V9}{01V9}], its generic points therefore lie in the smooth locus of \((X_{S'})_{\mathrm{red}}\to S'\). After shrinking the base once more to remove non-dominant components, this smooth locus contains a dense open of \((X_{S'})_{\mathrm{red}}\). Since \(F\) is constructible, there is a smaller dense open \(V\) on which the pullback of \(F\) is locally constant. Its complement contains no generic point of a component of the generic fibre. Hence the generic fibre of the complement has dimension less than \(n\).
\end{proof}

% object: theorem-generic-constructibility-base-change-rj (blueprint/chap-Constructibility-of-f/sec-Generic-base-change/theorem-generic-constructibility-base-change-rj/_section.md)
\phantomsection\label{bp51}%
\begin{bpstatementit}{Theorem~6.17}{ (Generic constructibility and base change for $Rj_*$)}
Let \(S\) be integral, and let \(j:X\hookrightarrow Y\) be an open immersion of \(S\)-schemes. For a constructible \(\Fell\)-sheaf \(F\) on \(X_{\acute{e}t}\), there is a nonempty open \(S_0\subset S\) such that

\[
Rj_{S_0,*}(F|_{X_{S_0}})\in\Dbc(Y_{S_0},\Fell),
\tag{6.19}
\]

and such that for every \(k\)-scheme \(T\to S_0\) of finite type, with base change \(j_T:X_T\hookrightarrow Y_T\) and projections \(g_X:X_T\to X_{S_0}\) and \(g_Y:Y_T\to Y_{S_0}\), the base change map

\[
\beta_T:g_Y^*Rj_{S_0,*}(F|_{X_{S_0}})
\longrightarrow Rj_{T,*}g_X^*(F|_{X_{S_0}})
\tag{6.20}
\]

is an isomorphism.
\end{bpstatementit}

% proof: proof/theorem-generic-constructibility-base-change-rj
\begin{proof}[{Proof}]\phantomsection\label{bp52}%
We work with the étale sheaf corresponding to \(F\), in \(D^+(-_{\acute{e}t},\Fell)\), and prove (6.19) and \(\beta_T\) together by induction on \(n=\dim X_{k(S)}\), allowing at each stage all integral bases \(S\) of finite type over \(k\). We first handle the complement of a finite set by passing to the base \(\mathbf A^1_S\), then the smooth, locally constant case, and finally the general case by dévissage.

\emph{Preliminary reductions.} Several modifications of the data are harmless, and we will make them freely in the induction. First, the assertions extend from open to locally closed immersions with the same source. Indeed, such an immersion factors as an open immersion into its scheme-theoretic closure followed by a closed immersion. Since closed direct image is exact, preserves constructibility, and commutes with arbitrary pullback [Stacks, \href{https://stacks.math.columbia.edu/tag/0959}{0959}], the assertions pass along this factorization. We will use the induction hypothesis in this form. Next, replacing \(Y\) by the scheme-theoretic closure of \(X\), we may assume that \(X\) is dense in \(Y\). We may also shrink \(S\) to be affine, and we may remove the closures of the images of the irreducible components of \(Y\) that do not dominate \(S\), since there are finitely many of them and none contains the generic point of \(S\). Finally, the assertions are local on \(Y\). Since we only use finite affine covers, the finitely many resulting opens of \(S\) can be intersected.

\emph{Base case.} If the generic fibre is empty or zero-dimensional, then after shrinking \(S\) the immersion \(j\) becomes an isomorphism or \(Y\) becomes empty. Suppose first that \(X_{k(S)}\) is empty. By density \(Y_{k(S)}\) is empty. On a finite affine cover, a coordinate algebra that vanishes over \(k(S)\) has its unit killed by a nonzero element of \(\mathcal O(S)\). Hence \(Y\) becomes empty after shrinking \(S\). If \(n=0\), the dense open \(X_{k(S)}\) of the zero-dimensional Noetherian space \(Y_{k(S)}\) is all of \(Y_{k(S)}\). The same argument applied to \(Y\setminus X\) shows that \(j\) becomes an isomorphism after shrinking \(S\).

Now let \(n\geq1\), and assume the theorem, including \(\beta_T\), for generic fibres of dimension less than \(n\).

\emph{Step 1: for \(Y\) projective over \(S\), both assertions hold on an open \(G\subset Y\) whose complement is finite over \(S\).} We replace the base \(S\) by the larger base \(\mathbf A^1_S\), over which the generic fibre is smaller, so that the induction hypothesis applies. To do this, choose charts and coordinates as in Lemma \hyperref[bp41]{6.13}. Via \(t_{\alpha r}\), we regard \(X\cap Y_\alpha\hookrightarrow Y_\alpha\) as an open immersion of schemes over the integral base \(\mathbf A^1_S\), which is of finite type over \(k\). By part (1) of that lemma, its generic fibre has dimension less than \(n\). Therefore the induction hypothesis, applied over the base \(\mathbf A^1_S\) to the restriction of \(F\), gives a dense open \(V_{\alpha r}\subset\mathbf A^1_S\) over which both assertions hold.

Since \(Rj_*\) commutes with restriction to the open \(Y_\alpha\), the assertions hold for \(Rj_*F\) on \(t_{\alpha r}^{-1}(V_{\alpha r})\). They continue to hold after a base change \(T\to S\) with \(T\) of finite type over \(k\), as we see by applying the hypothesis to the inverse image of \(V_{\alpha r}\) in \(\mathbf A^1_T\). Let \(G\) be the finite union of the opens \(t_{\alpha r}^{-1}(V_{\alpha r})\). Then on \(G\) the map \(\beta_T\) is an isomorphism, and the complex \(Rj_*F\) is bounded constructible. For the boundedness, note that there are finitely many pieces, and therefore the cohomology lies in a common finite range of degrees. Finally, part (2) of the lemma shows that, after shrinking \(S\), the reduced complement \(Z=Y\setminus G\) is finite over \(S\).

\emph{Step 2: the theorem holds in dimension \(n\) when \(X\) is smooth over \(S\) and \(F\) is locally constant.} We may assume that \(Y\) is projective over \(S\). Indeed, let \(Y_\beta\) be a member of a finite affine cover of \(Y\), embed it in some \(\mathbf A^m_S\) over the affine base, and let \(\overline Y\) be the closure of \(X\cap Y_\beta\) in \(\mathbf P^m_S\). Then \(X\cap Y_\beta\) is open and dense in \(\overline Y\), which is projective over \(S\). Moreover \(X\cap Y_\beta\) is still smooth over \(S\), with generic fibre of dimension at most \(n\). The result for \(X\cap Y_\beta\subset\overline Y\) gives it on \(Y_\beta\), because \(Rj_*\) and its base change maps commute with restriction to opens of the target.

For \(Y\) projective, Step 1 gives an open \(G\subset Y\) with complement finite over \(S\) on which both assertions hold. Since \(X\) is smooth over \(S\) and \(F\) is locally constant, Lemma \hyperref[bp45]{6.14} applies. It shows that the assertions hold on all of \(Y\) after shrinking \(S\).

\emph{Step 3: the general case in dimension \(n\).} We find a dense open \(V\subset X\) to which Step 2 applies, and treat the difference between \(F\) and \(Rv_*v^*F\), which lives on the smaller set \(X\setminus V\), by the induction hypothesis. Let \(S'\to S\) and \(V\) be as in Lemma \hyperref[bp49]{6.16}. Lemma \hyperref[bp47]{6.15} shows that it suffices to prove the theorem for \((X_{S'})_{\mathrm{red}}\hookrightarrow(Y_{S'})_{\mathrm{red}}\) and the pullback of \(F\) over \(S'\). Its generic fibre still has dimension \(n\) because \(S'\to S\) is a universal homeomorphism. We therefore replace \(S\) by \(S'\) and all schemes by these reductions. Then \(V\subset X\) is a dense open, smooth over \(S\), on which \(F\) is locally constant. The generic fibre of \(X\setminus V\) has dimension less than \(n\).

Let \(v:V\hookrightarrow X\) be the open immersion and \(i:X\setminus V\hookrightarrow X\) its reduced closed complement. By Step 2, after shrinking the base, both \(v\) and \(jv\) have bounded constructible direct images of \(v^*F\). Their base change maps are isomorphisms after every base change \(T\to S\) with \(T\) of finite type over \(k\). Complete the unit of \(v^*\dashv Rv_*\) to a triangle

\[
\Delta_F:\ F\longrightarrow Rv_*v^*F\longrightarrow D
\longrightarrow F[1].
\tag{6.21}
\]

The first term is the constructible sheaf \(F\), and the second is bounded constructible by Step 2. Hence \(D\) is bounded constructible too. Indeed, by the long exact cohomology sequence, its cohomology sheaves are extensions of kernels by cokernels of maps between constructible sheaves, and these are constructible [Stacks, \href{https://stacks.math.columbia.edu/tag/09BH}{09BH}]. Moreover \(v^*D=0\), as the unit is an isomorphism on \(V\). The localization triangle (6.18) for \(v\) and \(i\) now shows that

\[
D=i_*i^*D.
\]

The locally closed immersion \(ji:X\setminus V\to Y\) has generic fibre of dimension less than \(n\). We apply the induction hypothesis to each of the finitely many constructible cohomology sheaves of \(i^*D\), together with their truncation triangles. This gives

\[
R(ji)_*i^*D=Rj_*D\in\Dbc(Y,\Fell),
\]

with base change maps that are isomorphisms. Now apply \(Rj_*\) to the triangle \(\Delta_F\). In the resulting triangle, the middle term \(R(jv)_*v^*F\) is bounded constructible by Step 2, and the third term \(Rj_*D\) was treated just now. Hence \(Rj_*F\) is bounded constructible.

To see that the base change map \(\beta_T\) is an isomorphism, we compare the triangle \(\Delta_F\) with its analogue over \(T\). Let \(g_X:X_T\to X\), \(g_Y:Y_T\to Y\) and \(g_V:V_T\to V\) be the projections, and \(v_T:V_T\to X_T\) the base change of \(v\). The base change map

\[
\beta_v:g_X^*Rv_*v^*F\longrightarrow Rv_{T,*}g_V^*v^*F
\]

is an isomorphism by Step 2, and it carries \(g_X^*\) of the first arrow of the triangle \(\Delta_F\) to the unit \(g_X^*F\to Rv_{T,*}v_T^*g_X^*F\) of \(v_T^*\dashv Rv_{T,*}\). Since cones are not functorial, we have to choose the triangles compatibly, and we take as triangle for this unit the pullback of the triangle \(\Delta_F\), transported through \(\beta_v\). Now the base change map

\[
\beta_j:g_Y^*Rj_*\longrightarrow Rj_{T,*}g_X^*
\]

is a natural transformation of exact functors. It therefore induces a morphism from \(g_Y^*\) of \(Rj_*\) of the triangle \(\Delta_F\) to \(Rj_{T,*}\) of this triangle. On the second term it is \(\beta_j\) at \(Rv_*v^*F\) followed by \(Rj_{T,*}\beta_v\). Since base change maps compose [Stacks, \href{https://stacks.math.columbia.edu/tag/0EZR}{0EZR}, \href{https://stacks.math.columbia.edu/tag/0EZS}{0EZS}], this composite is the base change map for \(jv\), an isomorphism by Step 2. On the third term it is \(\beta_j\) at \(D\). It is an isomorphism because, composed with \(Rj_{T,*}\) of the base change isomorphism for the closed immersion \(i\) [Stacks, \href{https://stacks.math.columbia.edu/tag/0959}{0959}], it is the base change map for \(ji\) at \(i^*D\) treated above. Hence, by the five lemma for triangles, it is an isomorphism on the first term, which is \(\beta_T\).

This completes the induction.
\end{proof}

% text: prose-06-16
\phantomsection\label{bp53}%
The only nonempty open of \(\operatorname{Spec}k\) is \(\operatorname{Spec}k\) itself, and hence over a field the theorem holds without any shrinking.

% object: corollary-open-direct-image-over-field (blueprint/chap-Constructibility-of-f/sec-Generic-base-change/corollary-open-direct-image-over-field/_section.md)
\phantomsection\label{bp54}%
\begin{bpstatementit}{Corollary~6.18}{ (Open direct image over a field)}
Let \(j:X\hookrightarrow Y\) be an open immersion. Then

\[
Rj_*:\Dbc(X,A)\longrightarrow\Dbc(Y,A),
\tag{6.22}
\]

and for every \(k\)-scheme \(T\) of finite type, with projections \(g_X:X_T\to X\) and \(g_Y:Y_T\to Y\), and every \(K\in\Dbc(X,A)\), the base change map \(g_Y^*Rj_*K\to Rj_{T,*}g_X^*K\) is an isomorphism.
\end{bpstatementit}

% proof: proof/corollary-open-direct-image-over-field
\begin{proof}[{Proof}]\phantomsection\label{bp55}%
We argue on the étale site. The theorem gives the case of \(\Fell\)-sheaves, and we pass to \(\mathbf Z/\ell^r\)-sheaves by filtrations, to \(A\)-sheaves by forgetting the \(A\)-structure, and to complexes by truncation.

For a constructible \(\Fell\)-sheaf this is Theorem \hyperref[bp51]{6.17} with \(S=\operatorname{Spec}k\). A constructible sheaf \(F\) killed by \(\ell^r\) has the finite filtration \(F\supset\ell F\supset\cdots\supset\ell^rF=0\). Its successive quotients \(\ell^iF/\ell^{i+1}F\) are constructible and killed by \(\ell\), that is, constructible \(\Fell\)-sheaves. Each short exact sequence in the filtration is a triangle in the derived category, and \(Rj_*\) carries it to a triangle. By induction on \(r\), the long exact cohomology sequences of these triangles give (6.22). The base change maps form morphisms between the resulting triangles, and they are isomorphisms by the five lemma.

For a finite \(\ell\)-primary ring \(A\), we apply this to the underlying \(\mathbf Z/\ell^r\mathbf Z\)-sheaf and carry the \(A\)-action along all maps. This is possible for two reasons. First, \(Rj_*:D^+(X_{\acute{e}t},A)\to D^+(Y_{\acute{e}t},A)\) commutes with the forgetful functors to \(D^+(-_{\acute{e}t},\mathbf Z/\ell^r\mathbf Z)\) [Stacks, \href{https://stacks.math.columbia.edu/tag/0D6J}{0D6J}]. Second, an \(A\)-sheaf is constructible if and only if its underlying sheaf is. To see this, consider a stratum where the underlying sheaf is finite locally constant. There the action of each of the finitely many elements of \(A\) is a section of the finite locally constant sheaf of endomorphisms, hence constant étale locally. Therefore the \(A\)-module is locally constant on that stratum. Truncation triangles extend both conclusions from sheaves to bounded constructible complexes.
\end{proof}


% object: sec-Finite-coefficients-2 (blueprint/chap-Constructibility-of-f/sec-Finite-coefficients-2/_section.md), content hidden
\subsection*{\texorpdfstring{6.6\quad Finite coefficients}{6.6 Finite coefficients}}
\phantomsection\label{bp56}%
\addcontentsline{toc}{subsection}{\texorpdfstring{6.6\quad Finite coefficients}{6.6 Finite coefficients}}
% text: prose-06-17
\phantomsection\label{bp57}%
Let \(A\) be a finite commutative \(\ell\)-primary ring and \(f\) a separated morphism of finite type. We show that the right adjoint \(f^!\) of Theorem \bpout{5.31} preserves \(D^b_{\mathrm{ctf}}(-,A)\) (Theorem \hyperref[bp61]{6.20}). By the reduction at the start of this chapter and by the previous section, we only have to show that direct image along an open immersion \(j\) preserves finite Tor amplitude. The upper end of the Tor amplitude grows by at most the \(\ell\)-cohomological dimension \(\operatorname{cd}_\ell(j)\) of Definition \bpout{2.45}. This dimension is finite for morphisms of finite type over \(k\).

Finite Tor amplitude cannot be checked on cohomology sheaves, since these may have infinite Tor dimension over \(A\): for \(A=\mathbf Z/\ell^2\), already the constant sheaf \(\mathbf F_\ell\) has infinite Tor dimension. We therefore obtain the Tor bounds below from flat models. By [Stacks, \href{https://stacks.math.columbia.edu/tag/0F4M}{0F4M}], every object of \(D^b_{\mathrm{ctf}}(X,A)\) is represented by a bounded complex of constructible flat \(A\)-sheaves. On the strata of a finite stratification, such a model is locally a bounded complex of finite projective \(A\)-modules.

% object: proposition-open-direct-image-preserves-finite (blueprint/chap-Constructibility-of-f/sec-Finite-coefficients-2/proposition-open-direct-image-preserves-finite/_section.md)
\phantomsection\label{bp58}%
\begin{bpstatementit}{Proposition~6.19}{ (Open direct image preserves finite Tor amplitude)}
Let \(j:U\hookrightarrow X\) be an open immersion. Then \(Rj_*:D(U_{\mathrm{pro}\acute{e}t},A)\to D(X_{\mathrm{pro}\acute{e}t},A)\) restricts to

\[
Rj_*:D^b_{\mathrm{ctf}}(U,A)\longrightarrow D^b_{\mathrm{ctf}}(X,A).
\]

More precisely, if \(K\) has Tor amplitude in \([a,b]\) and \(d\ge
\operatorname{cd}_\ell(j)\), then \(Rj_*K\) has Tor amplitude in \([a,b+d]\).
\end{bpstatementit}

% proof: proof/proposition-open-direct-image-preserves-finite
\begin{proof}[{Proof}]\phantomsection\label{bp59}%
We test Tor amplitude by tensoring with modules. By a projection formula for constant modules we may move the tensor product inside \(Rj_*\). There the complex \(K\otimes^{\mathbf L}\underline M\) has amplitude \([a,b]\), and \(Rj_*\) adds at most \(d\) degrees.

The complex \(Rj_*K\) is bounded constructible by Corollary \hyperref[bp54]{6.18}, and \(\operatorname{cd}_\ell(j)\) is finite by Theorem theorem-results-from-etale-theory. To bound the Tor amplitude, we work on the étale site, where geometric points are conservative. By Corollary \bpout{2.44}, we have \(K\simeq\nu_U^*K_{\acute{e}t}\) for some \(K_{\acute{e}t}\in D^b_{\mathrm{ctf}}(U_{\acute{e}t},A)\). By Proposition \bpout{2.49}, the pro-étale pullback of \(Rj_{\acute{e}t,*}K_{\acute{e}t}\) is \(Rj_*K\).

We first prove the projection formula for constant coefficients. Namely, for every \(A\)-module \(M\), viewed as a constant sheaf, the projection morphism

\[
(Rj_*K)\otimes_A^{\mathbf L}\underline M_X
\xrightarrow{\ \sim\ }
Rj_*\bigl(K\otimes_A^{\mathbf L}\underline M_U\bigr)
\tag{6.23}
\]

is an isomorphism. Recall that for a scheme \(V\) with \(\operatorname{cd}_\ell(V)<\infty\), the projection formula of [Stacks, \href{https://stacks.math.columbia.edu/tag/0F12}{0F12}] identifies \(R\Gamma(V,K)\otimes^{\mathbf L}_AM\) with \(R\Gamma(V,K\otimes^{\mathbf L}_A\underline M)\). Moreover, a map of complexes of sheaves is an isomorphism if it induces isomorphisms on derived sections over a basis of the site. Since the sections of \(Rj_*\) over \(V\) are the sections over \(V\times_XU\), we apply the projection formula to every quasi-compact étale \(V\to X\) and to \(V\times_XU\). This gives isomorphisms between the derived sections of the two sides of the étale version of (6.23), compatibly with the projection morphism. Hence the étale version holds. Finally, inverse image is exact and symmetric monoidal, and it commutes with unbounded direct image, so the pullback of the étale isomorphism is (6.23).

On the étale site it suffices to bound Tor against constant \(A\)-modules. Indeed, let \(G\) be an étale sheaf of \(A\)-modules and \(\bar x\) a geometric point. The stalk of \(Rj_{\acute{e}t,*}K_{\acute{e}t}\otimes^{\mathbf L}_AG\) at \(\bar x\) is that of \(Rj_{\acute{e}t,*}K_{\acute{e}t}\otimes^{\mathbf L}_A\underline M\) with \(M=G_{\bar x}\), and geometric stalks are conservative. Now \(K\otimes_A^{\mathbf L}\underline M_U\) has cohomology in \([a,b]\). Hence the étale version of the right-hand side of (6.23) lies in \([a,b+d]\). By the projection formula, so does \(Rj_{\acute{e}t,*}K_{\acute{e}t}\otimes^{\mathbf L}_A\underline M\). In other words, \(Rj_{\acute{e}t,*}K_{\acute{e}t}\) has Tor amplitude in \([a,b+d]\).

To pass to \(X_{\mathrm{pro}\acute{e}t}\), we pull back a bounded flat constructible model of \(Rj_{\acute{e}t,*}K_{\acute{e}t}\) in degrees \([a,b+d]\) [Stacks, \href{https://stacks.math.columbia.edu/tag/0F4M}{0F4M}]. Since pullback preserves flatness, its pullback to \(X_{\mathrm{pro}\acute{e}t}\) is again a complex of flat sheaves in degrees \([a,b+d]\). Hence its tensor product with any sheaf of \(A\)-modules on \(X_{\mathrm{pro}\acute{e}t}\) is computed termwise and lies in degrees \([a,b+d]\). This argument also covers the sheaves on \(X_{\mathrm{pro}\acute{e}t}\) that are not pulled back from \(X_{\acute{e}t}\), for which geometric stalks would not detect the Tor amplitude.
\end{proof}

% text: prose-06-18
\phantomsection\label{bp60}%
For \(j:\mathbf G_m\hookrightarrow\mathbf A^1\) over \(\bar k\) and \(K=A\), the stalk of \(Rj_*A\) at the origin is the cohomology of the punctured strict henselization. It has cohomology \(A\) in degree \(0\) and \(A(-1)\) in degree \(1\), both free. It is therefore isomorphic to \(A\oplus A(-1)[-1]\). The Tor amplitude grows from \([0,0]\) to \([0,1]\), and the bound of the proposition, with \(d=\operatorname{cd}_\ell(j)=1\), is attained.

Combining the proposition with the computation of \(p^!\) for affine space, we obtain the main theorem for finite coefficients.

% object: theorem-f-shriek-finite-coefficients (blueprint/chap-Constructibility-of-f/sec-Finite-coefficients-2/theorem-f-shriek-finite-coefficients/_section.md)
\phantomsection\label{bp61}%
\begin{bpstatementit}{Theorem~6.20}{ ($f^!$ with finite coefficients)}
For every separated morphism \(f:X\to Y\) of finite type, the right adjoint \(f^!:\Loc(D^b_{\mathrm{ctf}}(Y,A))\to\Loc(D^b_{\mathrm{ctf}}(X,A))\) of Theorem \bpout{5.31} maps \(D^b_{\mathrm{ctf}}(Y,A)\) into \(D^b_{\mathrm{ctf}}(X,A)\). Consequently, its adjunction restricts to

\[
Rf_!:D^b_{\mathrm{ctf}}(X,A)\ \rightleftarrows\ D^b_{\mathrm{ctf}}(Y,A):f^!.
\]
\end{bpstatementit}

% proof: proof/theorem-f-shriek-finite-coefficients
\begin{proof}[{Proof}]\phantomsection\label{bp62}%
Locally on \(X\) and \(Y\), the morphism \(f\) is a closed immersion followed by the projection from an affine space. We therefore treat these two cases and then glue.

\emph{Closed immersions.} Let \(i:Z\hookrightarrow X\) be a closed immersion with open complement \(j:U\hookrightarrow X\), and let \(R\underline\Gamma_Z\) be the local cohomology functor of Section \bpout{2.14}. We show that \(R\underline\Gamma_Z\) preserves \(D^b_{\mathrm{ctf}}\) and that it is \(i^!\). For \(K\in D^b_{\mathrm{ctf}}(X,A)\), there is a triangle

\[
i_*R\underline\Gamma_ZK\longrightarrow K\longrightarrow Rj_*j^*K
\longrightarrow i_*R\underline\Gamma_ZK[1].
\]

By Proposition \hyperref[bp58]{6.19}, the two terms on the right lie in \(D^b_{\mathrm{ctf}}\), and hence so does \(i_*R\underline\Gamma_ZK\). Since \(i^*i_*=\mathrm{id}\) and \(i^*\) preserves \(D^b_{\mathrm{ctf}}\), we get \(R\underline\Gamma_ZK\in D^b_{\mathrm{ctf}}(Z,A)\). Explicitly, if \(K\) has Tor amplitude in \([a,b]\), then \(R\underline\Gamma_ZK\) has Tor amplitude in \([a,b+\operatorname{cd}_\ell(j)+1]\). The functor \(R\underline\Gamma_Z:D(X_{\mathrm{pro}\acute{e}t},A)\to D(Z_{\mathrm{pro}\acute{e}t},A)\) is right adjoint to \(i_*\) (Lemma lemma-sections-support-closed-subscheme), and \(\Loc(D^b_{\mathrm{ctf}}(-,A))\) is a full subcategory of \(D((-)_{\mathrm{pro}\acute{e}t},A)\). By uniqueness of right adjoints, \(R\underline\Gamma_ZK\) is therefore \(i^!K\) in \(\Loc(D^b_{\mathrm{ctf}}(Z,A))\).

\emph{Projections.} Let \(p:\mathbf A^d_Y\to Y\) be the projection. From the adjunction \(Rq_*\dashv q^*(1)[2]\) for the projective line \(q:\mathbf P^1_Y\to Y\) (Proposition \hyperref[bp22]{6.7}) and from Lemma \hyperref[bp26]{6.8}, we deduce that \(p^!\simeq p^*(-)(d)[2d]\) on \(\Loc(D^b_{\mathrm{ctf}}(Y,A))\). This functor preserves \(D^b_{\mathrm{ctf}}\), since pullback preserves a bounded constructible flat model and the Tate twist is locally free of rank one.

For example, for the origin \(i:0\hookrightarrow\mathbf A^1\) over \(\bar k\) and \(K=A\), the triangle gives \(i^!A\simeq A(-1)[-2]\) (Example \hyperref[bp79]{6.25} below). Its Tor amplitude \([2,2]\) lies in the predicted interval \([0,0+1+1]\).

\emph{Gluing.} After restricting to affine opens in \(X\) and \(Y\), a morphism of finite type factors as a closed immersion into an affine space over \(Y\) followed by its projection. Since \(f^!\) is compatible with composition and restriction to opens (Proposition \bpout{5.32}), the restriction of \(f^!K\) is obtained by applying the two functors just considered. Thus \(f^!K\) is bounded constructible of finite Tor amplitude locally on \(X\). A finite affine cover of the quasi-compact scheme \(X\) gives a single cohomological interval and a single interval of Tor amplitude valid everywhere. Finally, \(Rf_!\) preserves \(D^b_{\mathrm{ctf}}\) by Proposition \bpout{4.35}. Both sides of the adjunction therefore lie in \(D^b_{\mathrm{ctf}}\), and its Hom bijections restrict.
\end{proof}


% object: sec-ell-adic-coefficients-2 (blueprint/chap-Constructibility-of-f/sec-ell-adic-coefficients-2/_section.md), content hidden
\subsection*{\texorpdfstring{6.7\quad $\ell$-adic coefficients}{6.7 ell-adic coefficients}}
\phantomsection\label{bp63}%
\addcontentsline{toc}{subsection}{\texorpdfstring{6.7\quad $\ell$-adic coefficients}{6.7 ell-adic coefficients}}
% text: prose-06-19
\phantomsection\label{bp64}%
For \(\Lambda=\mathcal O_E\), \(E\) and \(\overline{\mathbf Q}_\ell\) we have \(D^b_{\mathrm{ctf}}=\Dbc\), and we show that \(f^!\) preserves \(\Dbc\). The argument of Section \hyperref[bp56]{6.6} uses two geometric facts, namely the finiteness of \(Rj_*\) for an open immersion and the adjunction for the projective line, and we deduce both from the case of torsion coefficients. This is possible because an integral complex lies in \(\Dbc\) if it is derived complete with constructible reduction modulo \(\pi_E\) (Proposition \bpout{3.11}), and rational complexes are built from integral ones (Lemma lemma-integral-models).

For \(Rj_*\) this reduction is immediate, since \(Rj_*\) commutes with \(R\varprojlim\) and with cones (Proposition \hyperref[bp65]{6.21}). For the projective line it is not, because the unit and counit are morphisms, and a compatible family of morphisms modulo \(\ell^n\) need not come from a morphism of the limits. We therefore construct the class of the diagonal and the trace once, at the level of inverse systems (Subsection \hyperref[bp67]{6.7.1}). Besides this, the integral case uses the structure of integral constructible complexes and the fact that proper direct image preserves direct sums (Subsection \bpout{5.4.2}).

% object: proposition-direct-image-along-open-immersion (blueprint/chap-Constructibility-of-f/sec-ell-adic-coefficients-2/proposition-direct-image-along-open-immersion/_section.md)
\phantomsection\label{bp65}%
\begin{bpstatementit}{Proposition~6.21}{ (Direct image along an open immersion)}
Let \(j:U\hookrightarrow X\) be a quasi-compact open immersion. Then \(Rj_*:D(U_{\mathrm{pro}\acute{e}t},\Lambda)\to D(X_{\mathrm{pro}\acute{e}t},\Lambda)\) maps \(\Dbc(U,\Lambda)\) into \(\Dbc(X,\Lambda)\) for \(\Lambda=\mathcal O_E\), \(E\) and \(\overline{\mathbf Q}_\ell\).
\end{bpstatementit}

% proof: proof/proposition-direct-image-along-open-immersion
\begin{proof}[{Proof}]\phantomsection\label{bp66}%
For \(\mathcal O_E\) we use the criterion of Proposition \bpout{3.11}, which asks for derived completeness and constructible reduction. Let \(K\in\Dbc(U,\mathcal O_E)\). Since \(Rj_*\) commutes with \(R\varprojlim\) [Stacks, \href{https://stacks.math.columbia.edu/tag/099L}{099L}], \(Rj_*K\) is derived complete. Since \(Rj_*\) also commutes with cones, the reduction of \(Rj_*K\) modulo \(\pi_E\) is \(Rj_*\) of the reduction of \(K\). That reduction comes from \(U_{\acute{e}t}\). Its direct image is therefore the pullback of the étale one [Stacks, \href{https://stacks.math.columbia.edu/tag/09A6}{09A6}], which is bounded constructible by Corollary \hyperref[bp54]{6.18}. Both conditions of the criterion hold, and so \(Rj_*K\in\Dbc(X,\mathcal O_E)\).

For \(E\) and \(\overline{\mathbf Q}_\ell\) we reduce to the integral case. Since \(Rj_*\) is exact, the objects \(K\) with \(Rj_*K\in\Dbc\) form a triangulated subcategory. It therefore suffices to treat the generators of Lemma lemma-integral-models. These come from integral objects by inverting \(\ell\) and by a colimit over finite extensions of \(E\). Both operations are filtered colimits of systems with a common lower bound, and \(Rj_*\) commutes with such colimits (Corollary \bpout{2.29}). Hence \(Rj_*\) of a generator is obtained from \(Rj_*\) of an integral object by the same operations, and it lies in \(\Dbc\).
\end{proof}

% object: subsec-The-projective-line (blueprint/chap-Constructibility-of-f/sec-ell-adic-coefficients-2/subsec-The-projective-line/_section.md), content hidden
\subsubsection*{\texorpdfstring{6.7.1\quad The projective line}{6.7.1 The projective line}}
\phantomsection\label{bp67}%
\addcontentsline{toc}{subsubsection}{\texorpdfstring{6.7.1\quad The projective line}{6.7.1 The projective line}}
% text: prose-06-20
\phantomsection\label{bp68}%
With torsion coefficients, \(Rq_*\dashv q^*(1)[2]\) is built from the cycle class \(\delta:\Delta_*\Lambda\to\Lambda(1)[2]\) of the diagonal of \(\mathbf P^1_Y\times_Y\mathbf P^1_Y\), the trace \(\operatorname{Tr}_q:Rq_*\Lambda(1)[2]\to\Lambda\) and two identities between them (Proposition \hyperref[bp22]{6.7}). The classes modulo \(\ell^n\), however, are morphisms in different derived categories, compatible under reduction. Since morphisms in a derived category are not computed levelwise, such a family need not give a morphism after \(R\varprojlim\). We therefore construct the class once, in the derived category of inverse systems of sheaves, and then pass to the limit.

% object: proposition-ell-adic-cycle-class-diagonal (blueprint/chap-Constructibility-of-f/sec-ell-adic-coefficients-2/subsec-The-projective-line/proposition-ell-adic-cycle-class-diagonal/_section.md)
\phantomsection\label{bp69}%
\begin{bpstatementit}{Proposition~6.22}{ (The $\ell$-adic cycle class of the diagonal)}
Let \(q:\mathbf P^1_Y\to Y\) be the projection, \(\Delta\) the diagonal of \(\mathbf P^1_Y\times_Y\mathbf P^1_Y\) and \(p_1,p_2\) its projections. There are a morphism \(\delta_{\mathcal O_E}:\Delta_*\mathcal O_E\to
\mathcal O_E(1)[2]\) reducing to the cycle classes of Lemma \hyperref[bp11]{6.3} and a trace \(\operatorname{Tr}_q:Rq_*\mathcal O_E(1)[2]\to\mathcal O_E\) lifting those of Lemma \hyperref[bp14]{6.4}, with \(\operatorname{Tr}_{p_i}\circ Rp_{i*}(\delta_{\mathcal O_E})=1\) for \(i=1,2\). Extending scalars gives \(\delta_\Lambda\) and \(\operatorname{Tr}_q\) with the same identities for \(\Lambda=E,\overline{\mathbf Q}_\ell\). All are compatible with base change along separated morphisms of finite type \(Y'\to Y\).
\end{bpstatementit}

% proof: proof/proposition-ell-adic-cycle-class-diagonal
\begin{proof}[{Proof}]\phantomsection\label{bp70}%
The cycle class comes from a single map defined with \(\mathbf Z\)-coefficients, the boundary of the divisor sequence. Hence all its reductions are reductions of one map of systems.

The diagonal is an effective Cartier divisor, which gives a boundary morphism \(\Delta_*\mathbf Z\to\mathbf G_m[1]\) (Lemma \hyperref[bp11]{6.3}). Tensoring with the system \((\mathbf Z/\ell^n)_n\) and using the Kummer sequence, we get one morphism of systems \(\Delta_*\mathbf Z/\ell^\bullet\to\mathbf Z/\ell^\bullet(1)[2]\). Its level \(n\) is the class of Lemma \hyperref[bp11]{6.3} for \(\mathbf Z/\ell^n\). The transition maps between levels are part of the data, and this is what a family of separately constructed classes would lack. Next we take \(R\varprojlim\). The site is replete with exact products (Propositions \bpout{2.22} and \bpout{Proposition 2.23}). Hence \(R\varprojlim\circ(-\otimes\mathbf Z/\ell^\bullet)\) is derived \(\ell\)-completion [Stacks, \href{https://stacks.math.columbia.edu/tag/099M}{099M}], and the limit \(\Delta_*\mathbf Z_\ell\to\mathbf Z_\ell(1)[2]\) is the completion of the boundary map of Lemma \hyperref[bp11]{6.3}. It is linear over the completed coefficient sheaf by [Stacks, \href{https://stacks.math.columbia.edu/tag/099I}{099I}]. We obtain the trace in the same way from the decomposition of \(Rq_*\mathbf Z/\ell^\bullet\).

Each \(\operatorname{Tr}_{p_i}\circ Rp_{i*}(\delta)\) is an endomorphism of \(\mathbf Z/\ell^\bullet\) in degree \(0\), that is, a family of endomorphisms of the levels. These are the identity by Proposition \hyperref[bp22]{6.7}. We then extend scalars to \(\mathcal O_E\), invert \(\ell\) and take the colimit over finite extensions (Corollary \bpout{2.29}). Base change holds because the Cartier and Kummer maps are natural and pullback commutes with \(R\varprojlim\) here (Lemma \bpout{2.54}).
\end{proof}

% text: prose-06-21
\phantomsection\label{bp71}%
To turn \(\delta_\Lambda\) and \(\operatorname{Tr}_q\) into an adjunction we repeat the torsion calculation. It used the projection formula and proper base change on the categories where \(Rq_*\) and \(q^*(1)[2]\) act. Proper base change on \(\Loc(D^b_{\mathrm{ctf}})\) is Proposition \bpout{5.24}, and we extend the projection formula in the same way.

% object: proposition-projection-formula-rational-coefficients (blueprint/chap-Constructibility-of-f/sec-ell-adic-coefficients-2/subsec-The-projective-line/proposition-projection-formula-rational-coefficients/_section.md)
\phantomsection\label{bp72}%
\begin{bpstatementit}{Proposition~6.23}{ (Projection formula with rational coefficients)}
For \(\Lambda=E\) or \(\overline{\mathbf Q}_\ell\), \(\Loc(D^b_{\mathrm{ctf}}(X,\Lambda))\) is closed under \(\otimes^{\mathbf L}\), and for proper \(p:X\to Y\) the map \((Rp_*M)\otimes^{\mathbf L}N\to Rp_*(M\otimes^{\mathbf L}p^*N)\) is an isomorphism for \(M\in\Loc(D^b_{\mathrm{ctf}}(X,\Lambda))\) and \(N\in\Loc(D^b_{\mathrm{ctf}}(Y,\Lambda))\).
\end{bpstatementit}

% proof: proof/proposition-projection-formula-rational-coefficients
\begin{proof}[{Proof}]\phantomsection\label{bp73}%
Both claims are first checked on \(\Dbc\) and then extended. Tensor products and both sides of the projection map commute with direct sums and cones in each variable. Hence, fixing all variables but one, the objects satisfying a claim form a localizing subcategory. It suffices that this subcategory contains \(\Dbc\) (Lemma \bpout{5.23}). By the dévissage of Corollary \bpout{3.20}, \(\Dbc\) is closed under \(\otimes^{\mathbf L}\), since \(a_!L\otimes^{\mathbf L}b_!M\) is the extension by zero of a lisse sheaf on the intersection of the two strata. The projection formula on \(\Dbc\) is Proposition \bpout{4.40}.
\end{proof}

% text: prose-06-22
\phantomsection\label{bp74}%
Now the torsion proof of the adjunction for \(\mathbf P^1\) can be run with \(\delta_\Lambda\) and \(\operatorname{Tr}_q\), once \(s_1=s_2=1\) is known.

% object: proposition-projective-line-rational-coefficients (blueprint/chap-Constructibility-of-f/sec-ell-adic-coefficients-2/subsec-The-projective-line/proposition-projective-line-rational-coefficients/_section.md)
\phantomsection\label{bp75}%
\begin{bpstatementit}{Proposition~6.24}{ (The projective line, rational coefficients)}
For \(\Lambda=E\) or \(\overline{\mathbf Q}_\ell\), the functor \(q^*(1)[2]:\Loc(D^b_{\mathrm{ctf}}(Y,\Lambda))\to\Loc(D^b_{\mathrm{ctf}}(\mathbf P^1_Y,\Lambda))\) is right adjoint to \(Rq_*\). Hence \(q^!\simeq q^*(1)[2]\) for the right adjoint \(q^!\) of Theorem \bpout{5.31}.
\end{bpstatementit}

% proof: proof/proposition-projective-line-rational-coefficients
\begin{proof}[{Proof}]\phantomsection\label{bp76}%
Define the unit and counit by the formulas (6.7) and (6.8) of Proposition \hyperref[bp22]{6.7}, using \(\Delta_*M\simeq p_1^*M\otimes^{\mathbf L}\Delta_*\Lambda\) [Stacks, \href{https://stacks.math.columbia.edu/tag/09AD}{09AD}] and the maps \(\delta_\Lambda\) and \(\operatorname{Tr}_q\) in place of their torsion versions. The proof of the identities \(s_1=s_2=1\) computed stalks of torsion sheaves, but apart from these identities the torsion calculation of the triangle identities used only the projection formula and proper base change. These hold here by Propositions \hyperref[bp72]{6.23} and \bpout{Proposition 5.24}. Thus the calculation reduces the triangle identities to the identities \(\operatorname{Tr}_{p_i}\circ Rp_{i*}(\delta_\Lambda)=1\) of Proposition \hyperref[bp69]{6.22}. Finally, uniqueness of adjoints gives the last claim.
\end{proof}


% object: subsec-Finiteness-of-f (blueprint/chap-Constructibility-of-f/sec-ell-adic-coefficients-2/subsec-Finiteness-of-f/_section.md), content hidden
\subsubsection*{\texorpdfstring{6.7.2\quad Finiteness of $f^!$}{6.7.2 Finiteness of f!}}
\phantomsection\label{bp77}%
\addcontentsline{toc}{subsubsection}{\texorpdfstring{6.7.2\quad Finiteness of $f^!$}{6.7.2 Finiteness of f!}}
% text: prose-06-23
\phantomsection\label{bp78}%
For admissible \(\Lambda\) and a closed immersion \(i:Z\hookrightarrow X\) with complement \(j\), the local cohomology functor \(R\underline\Gamma_Z\) is right adjoint to \(i_*=Ri_!\) (Lemma lemma-sections-support-closed-subscheme). Hence \(i^!\simeq R\underline\Gamma_Z\), and \(i^!\) is computed by the localization triangle

\[
i_*i^!K\longrightarrow K\longrightarrow Rj_*j^*K
\longrightarrow i_*i^!K[1].
\]

% object: example-closed-support (blueprint/chap-Constructibility-of-f/sec-ell-adic-coefficients-2/subsec-Finiteness-of-f/example-closed-support/_section.md)
\phantomsection\label{bp79}%
\begin{bpstatement}{Example~6.25}{ (Local cohomology at the origin of the affine line)}
Let \(X=\mathbf A^1_k\), \(k\) algebraically closed, \(Z\) the origin and \(U=\mathbf G_{m,k}\). Then \(Rj_*\Lambda_U\) has \(\mathcal H^0=\Lambda_X\) and \(\mathcal H^1=i_*\Lambda(-1)\), from the Kummer covers \(t\mapsto t^{\ell^n}\). The triangle gives \(i^!\Lambda_X\simeq\Lambda(-1)[-2]\), as Theorem \bpout{9.3}, applied to the structure maps of \(Z\) and \(X\), predicts.
\end{bpstatement}

% text: prose-06-24
\phantomsection\label{bp80}%
For \(\mathcal O_E\) we also need the projection formula on the localizing categories, which follows from the case of \(\Dbc\) and the results of Subsection \bpout{5.4.2}.

% object: lemma-projection-formula-integral-coefficients (blueprint/chap-Constructibility-of-f/sec-ell-adic-coefficients-2/subsec-Finiteness-of-f/lemma-projection-formula-integral-coefficients/_section.md)
\phantomsection\label{bp81}%
\begin{bpstatementit}{Lemma~6.26}{ (Projection formula with integral coefficients)}
For proper \(p:X\to Y\) the map \((Rp_*M)\otimes^{\mathbf L}N\to Rp_*(M\otimes^{\mathbf L}p^*N)\) is an isomorphism for \(M\in\Loc(D^b_{\mathrm{ctf}}(X,\mathcal O_E))\) and \(N\in\Loc(D^b_{\mathrm{ctf}}(Y,\mathcal O_E))\).
\end{bpstatementit}

% proof: proof/lemma-projection-formula-integral-coefficients
\begin{proof}[{Proof}]\phantomsection\label{bp82}%
On \(\Dbc\) this is Proposition \bpout{4.40}. Both sides make sense on the localizing categories, because \(M\otimes^{\mathbf L}p^*N\) lies in \(\Loc(D^b_{\mathrm{ctf}}(X,\mathcal O_E))\). For \(M\) and \(N\) in \(\Dbc\) this is Lemma \bpout{5.20}, and since the tensor product commutes with direct sums and cones in each variable, the argument of Proposition \hyperref[bp72]{6.23} extends it to the localizing categories. Indeed, fix one variable. Both sides are then exact and preserve direct sums in the other; for the variable \(M\) this holds because \(Rp_*\) preserves direct sums (Proposition \bpout{5.21}). Therefore Lemma \bpout{5.23}, applied first in \(M\) and then in \(N\), extends the isomorphism from \(\Dbc\) to the localizing categories.
\end{proof}

% text: prose-06-25
\phantomsection\label{bp83}%
We can now repeat the argument for finite coefficients with \(\mathcal O_E\), \(E\) and \(\overline{\mathbf Q}_\ell\).

% object: theorem-f-shriek-ell-adic-coefficients (blueprint/chap-Constructibility-of-f/sec-ell-adic-coefficients-2/subsec-Finiteness-of-f/theorem-f-shriek-ell-adic-coefficients/_section.md)
\phantomsection\label{bp84}%
\begin{bpstatementit}{Theorem~6.27}{ ($f^!$ with $\ell$-adic coefficients)}
Let \(\Lambda=\mathcal O_E\), \(E\) or \(\overline{\mathbf Q}_\ell\), and let \(f:X\to Y\) be separated of finite type. Then \(f^!\) maps \(D^b_{\mathrm{ctf}}(Y,\Lambda)\) into \(D^b_{\mathrm{ctf}}(X,\Lambda)\), and this restriction is right adjoint to \(Rf_!:D^b_{\mathrm{ctf}}(X,\Lambda)\to D^b_{\mathrm{ctf}}(Y,\Lambda)\).
\end{bpstatementit}

% proof: proof/theorem-f-shriek-ell-adic-coefficients
\begin{proof}[{Proof}]\phantomsection\label{bp85}%
As for finite coefficients, we factor \(f\) locally into a closed immersion and a projection from affine space, and treat the two factors separately.

Choose finite affine covers of \(X\) and \(Y\) and closed immersions \(U\hookrightarrow\mathbf A^d_V\). Here \(u:U\to X\) and \(v:V\to Y\) are the inclusions of the affine opens, and \(i:U\hookrightarrow\mathbf A^d_V\) and \(p:\mathbf A^d_V\to V\) are the two factors. Then \(u^*f^!\simeq i^!p^!v^*\) (Proposition \bpout{5.32}), and being in \(\Dbc\) is local on \(X\), so it suffices to show that \(p^!\) and \(i^!\) preserve \(\Dbc\).

For \(E\) and \(\overline{\mathbf Q}_\ell\) we argue as in Section \hyperref[bp56]{6.6}. For \(p\) we use Proposition \hyperref[bp75]{6.24}, the adjunction \(j_!\dashv j^*\) for \(\mathbf A^1_V\subset\mathbf P^1_V\), and Proposition \bpout{5.32} applied \(d\) times. They give \(p^!\simeq p^*(d)[2d]\), which preserves \(\Dbc\). For \(i\), let \(j'\) be the inclusion of the open complement of \(U\) in \(\mathbf A^d_V\). Then the localization triangle before Example \hyperref[bp79]{6.25}, with \(j'\) in place of \(j\), shows that \(i^!K\in\Dbc(U,\Lambda)\) for \(K\in\Dbc(\mathbf A^d_V,\Lambda)\). Indeed, its third term \(Rj'_*j'^*K\) is in \(\Dbc\) by Proposition \hyperref[bp65]{6.21}, and \(i^*i_*=\mathrm{id}\).

For \(\mathcal O_E\) that argument applies once its ingredients are available with integral coefficients. They are as follows. First, the adjunction for the projective line comes from Proposition \hyperref[bp69]{6.22}, with the triangle identities verified as in Proposition \hyperref[bp75]{6.24}. Next, the projection formula is Lemma \hyperref[bp81]{6.26}, and inverse image and tensor product preserve \(\Dbc(-,\mathcal O_E)\) by Lemma \bpout{5.20}. Moreover, the finiteness of \(Rj_*\) and proper base change are Propositions \hyperref[bp65]{6.21} and \bpout{Proposition 5.24}. We also use that derived completeness is local on \(X\); it is part of the criterion for \(\Dbc\) with \(\mathcal O_E\)-coefficients. Finally, \(Rf_!\) preserves \(\Dbc\) (Proposition \bpout{5.26}), and so the adjunction restricts to \(\Dbc\).
\end{proof}



% object: sec-Proof-of-the-theorem (blueprint/chap-Constructibility-of-f/sec-Proof-of-the-theorem/_section.md), content hidden
\subsection*{\texorpdfstring{6.8\quad Proof of the theorem}{6.8 Proof of the theorem}}
\phantomsection\label{bp86}%
\addcontentsline{toc}{subsection}{\texorpdfstring{6.8\quad Proof of the theorem}{6.8 Proof of the theorem}}
% text: prose-06-26
\phantomsection\label{bp87}%
The parts of Theorem \bpout{5.2} were proved separately for each kind of coefficient ring, and we collect the references.

% proof: proof/theorem-exceptional-inverse-image
\begin{proof}[{Proof of Theorem 5.2}]\phantomsection\label{bp88}%
Part (1) is Lemma \bpout{5.30}. Its compactness statement was proved separately for each kind of coefficient ring: for finite \(A\) in Proposition \bpout{5.14}, for \(E\) and \(\overline{\mathbf Q}_\ell\) in Theorem \bpout{5.13}, and for \(\mathcal O_E\) in Proposition \bpout{5.21}.

Part (2) is Proposition \bpout{5.26}. It rests on the fact that proper direct image preserves direct sums, which was proved for finite \(A\) in Corollary \bpout{5.16}, for \(E\) and \(\overline{\mathbf Q}_\ell\) in Lemma \bpout{5.18}, and for \(\mathcal O_E\) in Proposition \bpout{5.21} (see Proposition \bpout{5.22}).

Part (3) is Theorem \bpout{5.31}.

Part (4) is Theorem \hyperref[bp61]{6.20} for finite \(A\) and Theorem \hyperref[bp84]{6.27} for \(\mathcal O_E\), \(E\) and \(\overline{\mathbf Q}_\ell\). The second theorem includes \(\mathbf Z_\ell\) and \(\mathbf Q_\ell\) as the case \(E=\mathbf Q_\ell\); recall that \(D^b_{\mathrm{ctf}}=\Dbc\) in the integral and rational cases. For the last assertion, note that the categories \(D^b_{\mathrm{ctf}}(X,\Lambda)\) and \(D^b_{\mathrm{ctf}}(Y,\Lambda)\) are full subcategories of the localizing categories, and \(Rf_!\) and \(f^!\) preserve them. Hence the Hom bijections of the adjunction restrict to them.
\end{proof}



\clearpage\phantomsection\hypertarget{bp-outside}{}%
\section*{Outside this export}
\addcontentsline{toc}{section}{Outside this export}
References marked \ensuremath{{}^\ast} are to parts of the blueprint this document leaves out:
\begin{description}
\item[Proposition 2.22] Repleteness of the pro-étale topos, in Ⅰ The six operations (\bpslug{proposition-{}\allowbreak{}repleteness-{}\allowbreak{}pro-{}\allowbreak{}etale-{}\allowbreak{}topos})
\item[Proposition 2.23] Exact products, in Ⅰ The six operations (\bpslug{proposition-{}\allowbreak{}exact-{}\allowbreak{}products})
\item[Corollary 2.29] Continuity of cohomology, in Ⅰ The six operations (\bpslug{corollary-{}\allowbreak{}continuity-{}\allowbreak{}cohomology})
\item[Theorem 2.35] Bounded-below cohomology comparison, in Ⅰ The six operations (\bpslug{theorem-{}\allowbreak{}bounded-{}\allowbreak{}below-{}\allowbreak{}cohomology-{}\allowbreak{}comparison})
\item[Section 2.9] Direct images, in Ⅰ The six operations (\bpslug{sec-{}\allowbreak{}Direct-{}\allowbreak{}images})
\item[Theorem 2.38] Comparison of direct images, in Ⅰ The six operations (\bpslug{theorem-{}\allowbreak{}comparison-{}\allowbreak{}direct-{}\allowbreak{}images})
\item[Lemma 2.39] Compatibilities of the comparison of direct images, in Ⅰ The six operations (\bpslug{lemma-{}\allowbreak{}comparison-{}\allowbreak{}direct-{}\allowbreak{}images-{}\allowbreak{}compatibilities})
\item[Proposition 2.42] Locally constant sheaves of finite presentation, in Ⅰ The six operations (\bpslug{proposition-{}\allowbreak{}locally-{}\allowbreak{}constant-{}\allowbreak{}sheaves-{}\allowbreak{}finite-{}\allowbreak{}presentation})
\item[Theorem 2.43] Comparison for constructible sheaves, in Ⅰ The six operations (\bpslug{theorem-{}\allowbreak{}comparison-{}\allowbreak{}constructible-{}\allowbreak{}sheaves})
\item[Corollary 2.44] Finite Tor amplitude under comparison, in Ⅰ The six operations (\bpslug{corollary-{}\allowbreak{}finite-{}\allowbreak{}tor-{}\allowbreak{}amplitude-{}\allowbreak{}under-{}\allowbreak{}comparison})
\item[Definition 2.45] Cohomological dimension, in Ⅰ The six operations (\bpslug{definition-{}\allowbreak{}cohomological-{}\allowbreak{}dimension})
\item[Proposition 2.49] Unbounded comparison for finite coefficients, in Ⅰ The six operations (\bpslug{proposition-{}\allowbreak{}unbounded-{}\allowbreak{}comparison-{}\allowbreak{}finite-{}\allowbreak{}coefficients})
\item[Lemma 2.54] Inverse image of the coefficient sheaves, in Ⅰ The six operations (\bpslug{lemma-{}\allowbreak{}inverse-{}\allowbreak{}image-{}\allowbreak{}coefficient-{}\allowbreak{}sheaves})
\item[Section 2.14] Extension by zero and completion, in Ⅰ The six operations (\bpslug{sec-{}\allowbreak{}Extension-{}\allowbreak{}by-{}\allowbreak{}zero-{}\allowbreak{}and-{}\allowbreak{}completion})
\item[Proposition 3.11] The integral constructibility criterion, in Ⅰ The six operations (\bpslug{proposition-{}\allowbreak{}integral-{}\allowbreak{}constructibility-{}\allowbreak{}criterion})
\item[Example 3.19] A running example, in Ⅰ The six operations (\bpslug{example-{}\allowbreak{}running-{}\allowbreak{}example})
\item[Corollary 3.20] Generators of $D^b_c$ with field coefficients, in Ⅰ The six operations (\bpslug{corollary-{}\allowbreak{}coefficient-{}\allowbreak{}heart-{}\allowbreak{}b})
\item[Lemma 4.5] Identity and pasting, in Ⅰ The six operations (\bpslug{lemma-{}\allowbreak{}identity-{}\allowbreak{}pasting})
\item[Theorem 4.33] Pro-étale proper base change, in Ⅰ The six operations (\bpslug{theorem-{}\allowbreak{}proper-{}\allowbreak{}base-{}\allowbreak{}change-{}\allowbreak{}classical-{}\allowbreak{}pro})
\item[Proposition 4.35] Proper direct image for finite coefficients, in Ⅰ The six operations (\bpslug{proposition-{}\allowbreak{}proper-{}\allowbreak{}direct-{}\allowbreak{}image-{}\allowbreak{}finite-{}\allowbreak{}coefficients})
\item[Proposition 4.40] Proper direct image of constructible complexes, in Ⅰ The six operations (\bpslug{proposition-{}\allowbreak{}proper-{}\allowbreak{}direct-{}\allowbreak{}image-{}\allowbreak{}constructible-{}\allowbreak{}complexes})
\item[Theorem 5.2] The exceptional inverse image, in Ⅰ The six operations (\bpslug{theorem-{}\allowbreak{}exceptional-{}\allowbreak{}inverse-{}\allowbreak{}image})
\item[Theorem 5.13] Compactness with rational coefficients, in Ⅰ The six operations (\bpslug{theorem-{}\allowbreak{}compactness-{}\allowbreak{}rational-{}\allowbreak{}coefficients})
\item[Proposition 5.14] Compactness for finite coefficients, in Ⅰ The six operations (\bpslug{proposition-{}\allowbreak{}compactness-{}\allowbreak{}finite-{}\allowbreak{}coefficients})
\item[Proposition 5.15] Comparison with the étale derived category, in Ⅰ The six operations (\bpslug{proposition-{}\allowbreak{}comparison-{}\allowbreak{}etale-{}\allowbreak{}derived-{}\allowbreak{}category})
\item[Corollary 5.16] Proper direct image on the localizing categories, in Ⅰ The six operations (\bpslug{corollary-{}\allowbreak{}proper-{}\allowbreak{}direct-{}\allowbreak{}image-{}\allowbreak{}finite-{}\allowbreak{}coefficients-{}\allowbreak{}localizing})
\item[Lemma 5.18] Proper direct image of unbounded complexes, in Ⅰ The six operations (\bpslug{lemma-{}\allowbreak{}proper-{}\allowbreak{}direct-{}\allowbreak{}image-{}\allowbreak{}unbounded-{}\allowbreak{}complexes})
\item[Subsection 5.4.2] Integral coefficients, in Ⅰ The six operations (\bpslug{subsec-{}\allowbreak{}Integral-{}\allowbreak{}coefficients})
\item[Lemma 5.20] Inverse image and tensor product, in Ⅰ The six operations (\bpslug{lemma-{}\allowbreak{}inverse-{}\allowbreak{}image-{}\allowbreak{}tensor-{}\allowbreak{}product})
\item[Proposition 5.21] Compactness and proper direct image, in Ⅰ The six operations (\bpslug{proposition-{}\allowbreak{}compactness-{}\allowbreak{}proper-{}\allowbreak{}direct-{}\allowbreak{}image})
\item[Proposition 5.22] Proper direct image and extension by zero, in Ⅰ The six operations (\bpslug{proposition-{}\allowbreak{}proper-{}\allowbreak{}direct-{}\allowbreak{}image-{}\allowbreak{}extension-{}\allowbreak{}zero})
\item[Lemma 5.23] Checking isomorphisms on constructible complexes, in Ⅰ The six operations (\bpslug{lemma-{}\allowbreak{}checking-{}\allowbreak{}isomorphisms-{}\allowbreak{}constructible-{}\allowbreak{}complexes})
\item[Proposition 5.24] Proper base change on the localizing categories, in Ⅰ The six operations (\bpslug{proposition-{}\allowbreak{}proper-{}\allowbreak{}base-{}\allowbreak{}change-{}\allowbreak{}localizing-{}\allowbreak{}categories})
\item[Proposition 5.26] Construction and independence of $Rf_!$, in Ⅰ The six operations (\bpslug{proposition-{}\allowbreak{}construction-{}\allowbreak{}independence-{}\allowbreak{}rf-{}\allowbreak{}shriek})
\item[Proposition 5.28] Composition and base change for $Rf_!$, in Ⅰ The six operations (\bpslug{proposition-{}\allowbreak{}composition-{}\allowbreak{}base-{}\allowbreak{}change-{}\allowbreak{}rf-{}\allowbreak{}shriek})
\item[Proposition 5.29] Forgetting supports and composition, in Ⅰ The six operations (\bpslug{proposition-{}\allowbreak{}forgetting-{}\allowbreak{}supports-{}\allowbreak{}composition})
\item[Lemma 5.30] Hypotheses of Brown representability, in Ⅰ The six operations (\bpslug{lemma-{}\allowbreak{}hypotheses-{}\allowbreak{}brown-{}\allowbreak{}representability})
\item[Theorem 5.31] Existence of the exceptional inverse image, in Ⅰ The six operations (\bpslug{theorem-{}\allowbreak{}existence-{}\allowbreak{}exceptional-{}\allowbreak{}inverse-{}\allowbreak{}image})
\item[Proposition 5.32] Composition, open restriction and compactifications for $f^!$, in Ⅰ The six operations (\bpslug{proposition-{}\allowbreak{}composition-{}\allowbreak{}open-{}\allowbreak{}restriction-{}\allowbreak{}compactifications-{}\allowbreak{}f})
\item[Chapter 9] Smooth purity, in Ⅱ Duality and the trace formula (\bpslug{chap-{}\allowbreak{}Smooth-{}\allowbreak{}purity})
\item[Theorem 9.3] Smooth purity for all admissible coefficients, in Ⅱ Duality and the trace formula (\bpslug{theorem-{}\allowbreak{}smooth-{}\allowbreak{}purity-{}\allowbreak{}all-{}\allowbreak{}admissible-{}\allowbreak{}coefficients})
\item[Section 9.6] Compatibilities, in Ⅱ Duality and the trace formula (\bpslug{sec-{}\allowbreak{}Compatibilities})
\item[Chapter 10] Verdier and Poincaré duality, in Ⅱ Duality and the trace formula (\bpslug{chap-{}\allowbreak{}Verdier-{}\allowbreak{}and-{}\allowbreak{}Poincare-{}\allowbreak{}duality})
\item[Lemma 11.7] External duality, in Ⅱ Duality and the trace formula (\bpslug{lemma-{}\allowbreak{}external-{}\allowbreak{}duality-{}\allowbreak{}trace-{}\allowbreak{}morphism})
\item[Corollary 12.31] Point counting and the zeta function, in Ⅱ Duality and the trace formula (\bpslug{corollary-{}\allowbreak{}point-{}\allowbreak{}counting-{}\allowbreak{}zeta-{}\allowbreak{}function})
\item[Proposition D.3] Smooth proper curves, torsion prime to the characteristic, in Appendices (\bpslug{proposition-{}\allowbreak{}smooth-{}\allowbreak{}proper-{}\allowbreak{}curves-{}\allowbreak{}torsion-{}\allowbreak{}prime})
\end{description}
\clearpage\phantomsection%
\section*{References}
\addcontentsline{toc}{section}{References}
\bpbibitem{BBD82}{bp90}{A. Beilinson, J. Bernstein, and P. Deligne. \emph{Faisceaux pervers}, volume 100 of \emph{Astérisque}. 1982.}

\bpbibitem{BS15}{bp91}{B. Bhatt and P. Scholze. The pro-étale topology for schemes. \emph{Astérisque}, 369:99–201, 2015. \href{https://people.mpim-bonn.mpg.de/scholze/proetale.pdf}{Author's manuscript}.}

\bpbibitem{Con07}{bp92}{B. Conrad. Deligne's notes on Nagata compactifications. \emph{Journal of the Ramanujan Mathematical Society}, 22:205–257, 2007. \href{https://math.stanford.edu/\~conrad/papers/nagatafinal.pdf}{Author's manuscript}.}

\bpbibitem{Del74}{bp93}{P. Deligne. La conjecture de Weil. I. \emph{Publ. Math. IHÉS}, 43:273–307, 1974.}

\bpbibitem{Del80}{bp94}{P. Deligne. La conjecture de Weil. II. \emph{Publ. Math. IHÉS}, 52:137–252, 1980.}

\bpbibitem{Den22}{bp95}{C. Deninger. There is no “Weil-”cohomology theory with real coefficients for arithmetic curves. \emph{Ann. Sc. Norm. Super. Pisa Cl. Sci.}, 2022. To appear. arXiv:\href{https://arxiv.org/abs/2204.02714}{2204.02714}.}

\bpbibitem{Deu41}{bp96}{M. Deuring. Die Typen der Multiplikatorenringe elliptischer Funktionenkörper. \emph{Abh. Math. Sem. Univ. Hamburg}, 14:197–272, 1941.}

\bpbibitem{Eke90}{bp97}{T. Ekedahl. On the adic formalism. In \emph{The Grothendieck Festschrift, Vol. II}, volume 87 of \emph{Progress in Mathematics}, pages 197–218. Birkhäuser, 1990.}

\bpbibitem{Has36}{bp98}{H. Hasse. Zur Theorie der abstrakten elliptischen Funktionenkörper. I, II, III. \emph{J. Reine Angew. Math.}, 175:55–62, 69–88, 193–208, 1936.}

\bpbibitem{Hum72}{bp99}{J. E. Humphreys. \emph{Introduction to Lie Algebras and Representation Theory}, volume 9 of \emph{Graduate Texts in Mathematics}. Springer, 1972.}

\bpbibitem{Ill94}{bp100}{L. Illusie. Autour du théorème de monodromie locale. \emph{Astérisque}, 223:9–57, 1994. §4. \href{https://www.numdam.org/item/AST_1994__223__9_0/}{Original article}.}

\bpbibitem{Kat88}{bp101}{N. M. Katz. Travaux de Laumon. \emph{Astérisque}, 161–162:105–132, 1988. Séminaire Bourbaki, exp. 691.}

\bpbibitem{KW01}{bp102}{R. Kiehl and R. Weissauer. \emph{Weil conjectures, perverse sheaves and \(\ell\)-adic Fourier transform}, volume 42 of \emph{Ergebnisse der Mathematik}. Springer, 2001.}

\bpbibitem{Lau87}{bp103}{G. Laumon. Transformation de Fourier, constantes d'équations fonctionnelles et conjecture de Weil. \emph{Publ. Math. IHÉS}, 65:131–210, 1987.}

\bpbibitem{Mil86a}{bp104}{J. S. Milne. \href{https://www.jmilne.org/math/xnotes/AVs.pdf}{Abelian Varieties}. In \emph{Arithmetic Geometry}. 1986. Corrected online version.}

\bpbibitem{Mil86b}{bp105}{J. S. Milne. \href{https://www.jmilne.org/math/xnotes/JVs.pdf}{Jacobian Varieties}. In \emph{Arithmetic Geometry}. 1986. Corrected online version.}

\bpbibitem{Ran39}{bp106}{R. A. Rankin. Contributions to the theory of Ramanujan's function \(\tau(n)\) and similar arithmetical functions. II. \emph{Proc. Cambridge Philos. Soc.}, 35:357–372, 1939.}

\bpbibitem{SGA4}{bp107}{M. Artin, A. Grothendieck, and J.-L. Verdier, editors. \emph{Théorie des topos et cohomologie étale des schémas}, volumes 269, 270, 305 of \emph{Lecture Notes in Mathematics}. Springer, 1972–1973. SGA 4.}

\bpbibitem{SGA4½}{bp89}{P. Deligne. \emph{Cohomologie étale}, volume 569 of \emph{Lecture Notes in Mathematics}. Springer, 1977. SGA 4½; \emph{Théorèmes de finitude}, §3, and \emph{Rapport sur la formule des traces}, §3, and \emph{Sommes trigonométriques}, §1.}

\bpbibitem{SGA5}{bp108}{A. Grothendieck, editor. \emph{Cohomologie \(\ell\)-adique et fonctions \(L\)}, volume 589 of \emph{Lecture Notes in Mathematics}. Springer, 1977. SGA 5.}

\bpbibitem{SGA7II}{bp109}{P. Deligne and N. Katz, editors. \emph{Groupes de monodromie en géométrie algébrique}, volume 340 of \emph{Lecture Notes in Mathematics}. Springer, 1973. SGA 7 II; Exposé XIII, (2.1.7).}

\bpbibitem{Stacks}{bp110}{The Stacks Project Authors. \href{https://stacks.math.columbia.edu}{The Stacks Project}. Individual results are cited by their stable tags in the text.}

\bpbibitem{Var05}{bp111}{Y. Varshavsky. Lefschetz–Verdier trace formula and a generalization of a theorem of Fujiwara. 2005. Version 2. arXiv:\href{https://arxiv.org/abs/math/0505564}{math/0505564}. \href{https://arxiv.org/abs/math/0505564v2}{Author's preprint}.}

\bpbibitem{Wei48a}{bp112}{A. Weil. \emph{Sur les courbes algébriques et les variétés qui s'en déduisent}. Hermann, Paris, 1948.}

\bpbibitem{Wei48b}{bp113}{A. Weil. \emph{Variétés abéliennes et courbes algébriques}. Hermann, Paris, 1948.}

\bpbibitem{Wei49}{bp114}{A. Weil. Numbers of solutions of equations in finite fields. \emph{Bull. Amer. Math. Soc.}, 55:497–508, 1949.}

\end{document}
